\chapter{物理仿真与动力学}
\label{ch:07}

\paragraph*{本章阅读锚点}
本章的锚点是物理一致性。原始 3DGS 只要求图像重建误差小；物理仿真方法把高斯接入力学、碰撞、流体或材料模型，使场景能够按物理规律变化。


把一只高斯泼溅重建的杯子从桌上推下去——它会碎成几片？杯中液体如何飞溅？碎片的运动轨迹是否符合牛顿力学？3D Gaussian Splatting 将场景编码为数十万个离散椭球体，每个椭球只有位置、颜色、不透明度和协方差，既没有质量、也没有速度、更没有内部应力。要让这些"没有灵魂的彩色气泡"遵循物理定律运动，研究者必须在纯视觉表示和物理引擎之间架起桥梁。这一根本矛盾催生了 3DGS 物理仿真领域：如何赋予高斯基元物理属性，使其在可微分渲染管线中产生物理可信的动态行为？

围绕这一核心挑战，研究从多个方向展开。第~\ref{sec:07_1}~节讨论\textbf{物理约束与仿真}——从物质点法（MPM）耦合到可微分物理引擎集成，从物理属性自动推断到端到端物理感知 4D 仿真，涵盖刚体、弹性体、流体和断裂等多种材料行为。该节进一步细分为五个子方向：物理仿真集成（紧耦合/松耦合/感知驱动三条路线）、流体与粒子动力学（流体模拟与重建）、可变形物体与布料仿真（弹性体/布料/4D 动态）、声音与多物理场（声场/火焰/降雨/风/物理场可视化）、以及其他物理方法（哈密顿力学/拓扑推断/VR 交互）。

\begin{figure}[htbp]
\centering
\resizebox{\textwidth}{!}{\begin{tikzpicture}[
  route/.style={rectangle, rounded corners=5pt, draw=#1, fill=#1!10,
    text width=3.2cm, minimum height=0.9cm, align=center, font=\small},
  method/.style={rectangle, rounded corners=3pt, draw=#1, fill=#1!6,
    font=\scriptsize, text width=3.0cm, align=center, minimum height=0.6cm},
  arr/.style={-{Stealth[length=4pt]}, techgray!60},
  >=Stealth
]

  \node[route=techgray, text width=4cm, font=\small\bfseries] (root) at (39.9, 5.6) {物理仿真与动力学方法分类};
  \node[route=techblue] (b0) at (10.0, 3.1) {物理仿真集成};
  \node[method=techblue] (b0_l0) at (1.9, 0.6) {紧耦合MPM\\{\tiny PhysGaussian, i-PhysGaussian, GaussianFluent, PhysMorph-GS}};
  \node[method=techblue] (b0_l1) at (6.2, 0.6) {松耦合代理\\{\tiny GASP, MaGS, PhysTalk}};
  \node[method=techblue] (b0_l2) at (10.0, 0.6) {属性感知\\{\tiny Sim Anything, PUGS, PhysGS, PIXIE, FastPhysGS}};
  \node[method=techblue] (b0_l3) at (14.2, 0.6) {连续介质/功能\\{\tiny PIDG, PG-3DGS, ReconPhys, PIG, Mirage2Matter}};
  \node[route=techgreen] (b1) at (25.1, 3.1) {流体与粒子};
  \node[method=techgreen] (b1_l0) at (18.9, 0.6) {无网格求解器\\{\tiny Gaussian Fluids, LagrangianSplats}};
  \node[method=techgreen] (b1_l1) at (23.0, 0.6) {流体重建\\{\tiny FluidGaussian, FLUIDSPLAT}};
  \node[method=techgreen] (b1_l2) at (27.1, 0.6) {流体资产提取\\{\tiny 3D Dynamic Fluid Assets}};
  \node[route=techorange] (b2) at (39.9, 3.1) {可变形体/布料};
  \node[method=techorange] (b2_l0) at (32.1, 0.6) {弹簧-质点/端到端\\{\tiny PhysTwin, NeuMA, 3DGSim}};
  \node[method=techorange] (b2_l1) at (37.4, 0.6) {混合动力学/单目\\{\tiny Particle-Grid, MonoPhysics, PhysConvex}};
  \node[method=techorange] (b2_l2) at (42.4, 0.6) {布料/细长几何\\{\tiny CloDS, Cloth-Splatting, SketchRodGS}};
  \node[method=techorange] (b2_l3) at (47.0, 0.6) {4D参数优化\\{\tiny PhysFlow, NGFF, PEGS}};
  \node[route=techpurple] (b3) at (53.8, 3.1) {声音与多物理场};
  \node[method=techpurple] (b3_l0) at (51.5, 0.6) {声场重建/合成\\{\tiny AudioGS, SonicGauss}};
  \node[method=techpurple] (b3_l1) at (56.1, 0.6) {火焰/降雨/风\\{\tiny FlameGS, FieryGS, RainyGS, DiffWind}};
  \node[method=techpurple] (b3_l2) at (60.6, 0.6) {物理场可视化\\{\tiny GS-Surrogate}};
  \node[route=techgray] (b4) at (66.9, 3.1) {其他物理方法};
  \node[method=techgray] (b4_l0) at (65.0, 0.6) {哈密顿力学\\{\tiny NeHaD}};
  \node[method=techgray] (b4_l1) at (68.8, 0.6) {拓扑推断\\{\tiny TopoGaussian, DynamicTree}};
  \node[method=techgray] (b4_l2) at (72.8, 0.6) {VR/系统集成\\{\tiny LIVE-GS}};
  \draw[arr] (root) -- (b0);
  \draw[arr] (root) -- (b1);
  \draw[arr] (root) -- (b2);
  \draw[arr] (root) -- (b3);
  \draw[arr] (root) -- (b4);
  \draw[arr] (b0) -- (b0_l0);
  \draw[arr] (b0) -- (b0_l1);
  \draw[arr] (b0) -- (b0_l2);
  \draw[arr] (b0) -- (b0_l3);
  \draw[arr] (b1) -- (b1_l0);
  \draw[arr] (b1) -- (b1_l1);
  \draw[arr] (b1) -- (b1_l2);
  \draw[arr] (b2) -- (b2_l0);
  \draw[arr] (b2) -- (b2_l1);
  \draw[arr] (b2) -- (b2_l2);
  \draw[arr] (b2) -- (b2_l3);
  \draw[arr] (b3) -- (b3_l0);
  \draw[arr] (b3) -- (b3_l1);
  \draw[arr] (b3) -- (b3_l2);
  \draw[arr] (b4) -- (b4_l0);
  \draw[arr] (b4) -- (b4_l1);
  \draw[arr] (b4) -- (b4_l2);
\end{tikzpicture}}
\caption{物理仿真与动力学方法分类}
\label{fig:ch07_physics_taxonomy}
\end{figure}

%=====================================================================
\section{物理约束与仿真}
\label{sec:07_1}

本节按材料类型和技术路线组织，涵盖物理仿真集成、流体与粒子动力学、可变形物体与布料仿真、声音与多物理场，以及其他物理方法共五个子方向。

\begin{figure}[htbp]
\centering
\resizebox{\textwidth}{!}{\begin{tikzpicture}[
  box/.style={rectangle, rounded corners=4pt, draw=#1, fill=#1!10,
    text width=3.2cm, align=center, font=\small, minimum height=1cm},
  sbox/.style={rectangle, rounded corners=3pt, draw=#1!60, fill=#1!5,
    font=\scriptsize, text width=2.4cm, align=center, minimum height=0.5cm},
  arr/.style={-{Stealth[length=6pt]}, thick, techgray},
  darr/.style={-{Stealth[length=6pt]}, thick, techgray, dashed},
  >=Stealth
]

  \node[box=techblue] (g1) at (3.9, 0.6) {3D高斯参数(位置/协方差/SH)};
  \node[box=techgreen] (g2) at (13.0, 0.6) {物理代理构建(网格/粒子/凸包)};
  \node[box=techorange] (g3) at (21.9, 0.6) {物理引擎(SPH/FEM/PBD)};
  \node[box=techpurple] (g4) at (30.3, 0.6) {约束求解(碰撞/形变/断裂)};
  \node[box=techblue] (g5) at (38.9, 0.6) {高斯参数更新(位置/旋转/形变)};
  \node[font=\scriptsize, text=techgray] at (13.0, 1.8) {黑盒引擎集成};
  \node[font=\scriptsize, text=techgray] at (21.9, -0.6) {材料属性估计};
  \draw[arr] (g1) -- (g2);
  \draw[arr] (g2) -- (g3);
  \draw[arr] (g3) -- (g4);
  \draw[arr] (g4) -- (g5);
  \draw[darr] (g5) to[bend right=30] node[midway, font=\scriptsize, text=techgray] {循环迭代} (g2);
\end{tikzpicture}}
\caption{物理仿真集成框架：从高斯参数到物理代理再到约束求解的闭环}
\label{fig:ch07_physics_sim}
\end{figure}

\subsection{物理仿真集成}

将物理仿真引入 3DGS 场景，首要问题是选择何种耦合方式——高斯基元如何与物理引擎对接，以及物理属性从何而来。这一领域的技术演进经历了三个阶段：从直接将高斯基元作为 MPM（Material Point Method，物质点法）粒子参与仿真的"紧耦合"路线，到通过代理结构桥接高斯与黑盒物理引擎的"松耦合"路线，再到利用大语言模型（LLM，Large Language Model）和视觉语言模型（VLM，Vision-Language Model）自动推断物理属性的"感知驱动"路线。

\begin{table}[htbp]
\centering
\caption{3DGS 物理仿真方法对比}
\label{tab:07_physics_methods}
\resizebox{\textwidth}{!}{
\begin{tabular}{lccccc}
\hline
\textbf{方法} & \textbf{耦合策略} & \textbf{物理引擎} & \textbf{材料类型} & \textbf{属性获取} & \textbf{速度} \\
\hline
PhysGaussian & 紧耦合（MPM） & 显式 MPM & 弹性/塑性 & 手动 & -- \\
i-PhysGaussian & 紧耦合（隐式 MPM） & 隐式 Newton-GMRES & 弹性/塑性 & 手动 & 20$\times$ 更大 $\Delta t$ \\
GaussianFluent & 紧耦合（CD-MPM） & CD-MPM & 断裂/流体 & 手动+生成 & 实时 \\
GASP & 松耦合 & 任意黑盒 & 多材料 & 黑盒引擎 & -- \\
Sim Anything & 感知驱动 & MPM & 多材料 & MLLM 零样本 & 2 分钟 \\
PhysFlow & 感知驱动 & 可微分 MPM & 7 种材料 & GPT-4 + 光流 & -- \\
NeuMA & 端到端 & Particle-GS & 弹性 & 物理先验+修正 & -- \\
\hline
\end{tabular}
}
\end{table}

物质点法（Material Point Method, MPM）是连接 3DGS 与物理仿真最自然的选择——高斯基元本身就是三维空间中的点，可以直接充当 MPM 粒子。MPM 的核心思想是在拉格朗日粒子和欧拉网格之间交替传递信息：粒子携带质量、速度和变形梯度等物理量，在每一步中将信息转移到背景网格上求解动量方程，再将更新后的速度和位置传回粒子。

\subsubsection{紧耦合路线：MPM直接物理仿真}

本路线介绍将高斯基元直接作为MPM粒子参与物理仿真的方法，核心优势是无需格式转换即可实现"所见即所仿"，但限制为MPM一种物理引擎。

\paragraph{PhysGaussian：统一粒子表示的紧耦合仿真}

PhysGaussian（2023.11）~\cite{2311_12198}修改了高斯基元与物理仿真的对接方式——用\textbf{统一的粒子表示}替代传统的\textbf{仿真-渲染分离管线}。在标准 3DGS 中，高斯基元仅是渲染基元，不携带任何物理属性。若要为场景生成物理动态，传统管线需要先将高斯转换为四面体网格或笼状代理（cage mesh），再在仿真器中驱动网格运动，最后将结果映射回渲染表示。这一多阶段流程引入了分辨率不匹配和仿真-渲染间的离散化误差。PhysGaussian 的核心洞察是：高斯基元本身就是三维空间中的点，可以直接充当 MPM 粒子，从而实现``所见即所仿''（what you see is what you simulate, WS$^2$）。

具体而言，PhysGaussian 为每个高斯基元赋予质量 $m_p$、初始体积 $V_p^0$、速度 $v_p$ 和\textbf{变形梯度} $F_p$（$F_p$ 描述该粒子从初始材料状态到当前世界状态的局部仿射变换，包含拉伸、旋转和剪切信息；可进一步分解为弹性部分和塑性部分 $F = F^E F^P$ 以支持永久形变）。MPM 仿真循环中，粒子先将质量和动量通过 B-spline 核函数 $w_{ip}$ 传递到背景欧拉网格，在网格上求解动量方程得到更新速度 $\mathbf{v}_i^{n+1}$，再传回粒子：

\begin{equation}
\frac{m_i}{\Delta t}(\mathbf{v}_i^{n+1} - \mathbf{v}_i^{n}) = -\sum_p V_p^0 \frac{\partial \Psi}{\partial \mathbf{F}}(\mathbf{F}_p^{E,n}) {\mathbf{F}_p^{E,n}}^\top \nabla w_{ip}^n + \mathbf{f}_i^{\text{ext}},
\end{equation}

其中 $\Psi$ 为超弹性势能密度，等号右侧求和项将各粒子的内力（应力）汇聚到网格节点，$\mathbf{f}_i^{\text{ext}}$ 为外力。网格速度回传后，粒子的位置和变形梯度同步更新：

\begin{equation}
\mathbf{x}_p^{n+1} = \mathbf{x}_p^n + \Delta t \, \mathbf{v}_p^{n+1}, \quad \mathbf{F}_p^{E,n+1} = (\mathbf{I} + \Delta t \, \nabla \mathbf{v}_p) \, \mathbf{F}_p^{E,n}.
\end{equation}

这意味着每一步仿真不仅移动了高斯的位置，还更新了其变形状态。

关键创新在于如何将变形梯度映射为高斯的协方差矩阵。直观地说，如果一个椭球在材料空间中是 $A_p$，经过局部仿射变换 $F_p$ 后，它在世界空间中的形状就变成 $F_p A_p F_p^\top$。就像捏住一块橡皮泥的两端拉伸，整个形状随之变形。PhysGaussian 利用这一关系直接推导出仿真后高斯的协方差：

\begin{equation}
\mathbf{a}_p(t) = \mathbf{F}_p(t) \, \mathbf{A}_p \, \mathbf{F}_p(t)^\top,
\end{equation}

其中 $\mathbf{A}_p$ 为初始材料空间中的协方差（即标准 3DGS 重建后的原始椭球形状），$\mathbf{a}_p(t)$ 为仿真 $t$ 时刻的世界空间协方差。该映射的关键优势在于：仿真的输出可以直接用于渲染，无需任何格式转换。MPM 粒子就是高斯基元，变形梯度就是协方差的更新规则。不透明度 $\alpha_p$ 和球谐系数 $C_p$ 在仿真过程中保持不变，视为材料的固有视觉属性。该方法支持弹性体、塑性金属、非牛顿流体和颗粒材料等多种本构模型，并通过回映映射（return mapping）$\mathcal{Z}$ 将超出弹性区域的变形梯度投影回屈服面内以处理塑性流动。

然而，标准 3DGS 重建只在高斯表面放置粒子，内部是空心的——一个空心杯子在物理仿真中会塌陷。因此 MPM 耦合方法都需要\textbf{内部粒子填充}（interior filling）：通过体素化分析判定内部区域，在新填充的体素中初始化高斯粒子，并从最近的表面高斯继承外观参数。这也是后续多个 MPM 方法共享的前置步骤。

\paragraph{i-PhysGaussian：隐式时间积分突破稳定性瓶颈}

从 PhysGaussian 出发，i-PhysGaussian（2026.02）~\cite{2602_17117}针对其显式时间积分的稳定性瓶颈提出根本改进。显式 MPM 使用步初量直接预测步末状态（前向 Euler），时间步 $\Delta t$ 受到 CFL（Courant-Friedrichs-Lewy，库朗-弗里德里希斯-列维）条件严格限制，超出阈值即产生数值发散。i-PhysGaussian 将时间积分替换为隐式 Newton-GMRES（广义最小残差法）求解器：在每一步中，不直接接受显式更新结果，而是求解一个动量残差方程
\begin{equation}
R_I(\mathbf{u}) = f_I^{\text{ext}}(\mathbf{u}) + f_I^{\text{int}}(\mathbf{u}) - m_I \, a_I^{n+1}(\mathbf{u}_I),
\end{equation}

其中 $\mathbf{u}$ 为网格位移增量，$f_I^{\text{int}}$ 和 $f_I^{\text{ext}}$ 分别为内力和外力，$a_I^{n+1}$ 为 Newmark 关系下的步末加速度。该残差通过 Newton 外循环（线性化残差）和 GMRES 内循环（矩阵无关的 Krylov 子空间求解）联合优化至收敛，使步末状态严格满足动量平衡。实际效果是：在高达 20 倍更大的时间步下仍保持稳定，且对材料刚度不敏感。这意味着同样的设置可以处理从软橡皮泥到硬金属等不同刚度的材料，无需针对每种材料反复调试时间步。

\paragraph{GaussianFluent：脆性断裂与混合材料仿真}

在 MPM 紧耦合路线的基础上，GaussianFluent（2026.01）~\cite{2601_09265}首次在 3DGS 框架中实现了\textbf{脆性断裂}和\textbf{混合材料}仿真。断裂仿真的核心挑战是标准 3DGS 缺乏内部结构——切开一个西瓜需要看到红色的瓤和黑色的籽。GaussianFluent 提出基于生成模型的无训练内部纹理合成管线：先将物体沿 X 轴均匀切片，对每片使用 MVInpainter 和 SD-XL 进行纹理修复，再沿三个正交轴迭代细化以确保三维一致性。

在物理引擎方面，GaussianFluent 引入连续损伤 MPM（CD-MPM，Continuum Damage MPM）替代标准 MPM：为每个粒子跟踪损伤变量 $\kappa$，当变形梯度的体积拉伸 $p$ 和剪切畸变 $q$ 超出屈服面 $y(p, q; p_0, \kappa, M) = 0$ 时，粒子的应力生成能力随 $\kappa$ 递增而衰减。关键创新是\textbf{连续回映映射}（continuous return mapping）：原始 CD-MPM 的屈服椭圆投影在 $p = p_0$ 边界处存在不连续跳跃，一个机器精度波动可能导致完全不同的投影结果。GaussianFluent 引入动态投影中心 $p_c' = p_c + \delta_k(p^{\text{tr}})(p^{\text{tr}} - p_c)$ 替代固定中心，当 $k=2$ 时实现光滑连续的映射。通过为物体的不同部位（如西瓜的籽、瓤、皮）赋予不同 $\kappa$ 值，GaussianFluent 实现了混合材料的差异化断裂行为。

\paragraph{PhysMorph-GS：目标驱动的可微分闭环变形}

沿同一紧耦合路线，PhysMorph-GS（2025.11）~\cite{2511_16988}修改了 PhysGaussian 的开环仿真范式——用目标驱动的可微分优化闭环替代手动设定物理参数后直接前向仿真。PhysGaussian 中，用户设定材料参数后仿真按自身规律运行，结果如何全凭物理决定。PhysMorph-GS 的目标则是指定一个目标形状（如从球体变形为牛），让物理仿真自动产生走向该目标的变形轨迹。其核心洞察是：渲染梯度应当通过变形梯度 $F$（而非粒子位置 $x$）注入仿真循环。直觉上，MPM 的驱动力来自应力，而应力由 $F$ 决定——直接在 $F$ 空间施加梯度与物理求解器天然对齐，无需跨空间（位置$\to$应力）的任意缩放。相比之下，$L/x$ 仅作用于表面附近粒子，且与物理步的应力更新缺乏量纲对应，容易引发数值不稳定。此外，PhysMorph-GS 引入 Chamfer 引导塑性（Chamfer-guided plasticity）：当弹性回复力威胁抵消已获得的变形时，通过乘法分解 $F = F_e F_p$ 逐步迁移材料静止状态，将物理从"阻力"转化为"盟友"。

\paragraph{Resonance4D：从观测视频反演全物理参数}

与 PhysMorph-GS 给定目标形状做主动变形的目标不同，Resonance4D（2026.04）~\cite{2604_01994}从观测视频中反推物理参数并生成动态响应。与 PhysGaussian 的"给定参数$\to$仿真$\to$渲染"单向管线相比，Resonance4D 实现了反向路径：从视觉观测反演物理参数。

该方法的关键瓶颈在于运动监督信号的质量。此前方法依赖在线视频扩散模型提供帧级运动引导，峰值 GPU 内存高达 35\,GB。Resonance4D 提出双域运动监督（Dual-domain Motion Supervision, DMS）替代这一重计算开销：空间域使用 SSIM 约束帧级结构一致性；频谱域对帧差序列执行分块 3D FFT（快速傅里叶变换），匹配局部频谱的幅度（$\lambda_{\text{mag}}$）和相位（$\lambda_{\text{phase}}$）。帧差设计抑制静态外观、突出运动模式，分块策略避免全局 FFT 将运动信号与背景混合。这一轻量级监督将峰值 GPU 内存从 35\,GB 降至约 20\,GB，首次实现全参数（杨氏模量、泊松比、屈服应力、塑性粘度、内摩擦角、密度共 6 项）的联合优化。Resonance4D 还引入基于拉丁超立方采样（LHS，Latin Hypercube Sampling）的仿真驱动初始化策略：在参数空间的对数域采样多组候选，各运行短仿真后按 MS-SSIM 选择最优起点，显著提升多参数联合优化的收敛稳定性。

\paragraph{PIDG：连续介质力学驱动的无引擎仿真}

走出与 PhysGaussian 完全不同的技术路线，PIDG（2025.11）~\cite{2511_06299}不使用 MPM 等显式物理引擎，而是从连续介质力学出发，将每个高斯粒子视为具有时变本构参数的拉格朗日物质点。PhysGaussian 通过 MPM 的粒子-网格传输求解运动方程，PIDG 则将物理约束直接编码为损失函数：Cauchy 动量残差 $r(x,t) = \rho(\frac{\partial v}{\partial t} + (v \cdot \nabla)v) - \nabla \cdot \sigma$ 衡量每个粒子的动量平衡偏离，其 $L_2$ 范数作为物理正则化项 $L_{\text{CMR}}$ 加入优化。

该方法的核心设计是时变材料场（time-evolving material field）：一个多感知机（MLP，Multi-Layer Perceptron）联合预测每个粒子的速度 $v$ 和本构应力 $\sigma$ 的六个独立分量。无需指定具体的本构模型（弹性/流体/刚体），网络自动学习适合当前场景的应力-应变关系。在场景表示层面，PIDG 采用静态-动态解耦的 4D 分解哈希编码（四个 3D 哈希网格分别覆盖 xyz、xyt、yzt、xzt），先联合优化静态几何与动态变形，再冻结静态部分仅精调动态区域。

\textbf{动机。}传统 4D 表示将静态与动态耦合在统一场中，静态结构的梯度被动态区域的优化噪声干扰，解耦编码使静态几何获得稳定的优化信号。在监督层面，PIDG 提出拉格朗日粒子流匹配：将高斯基元按速度场推进后投影到图像平面，生成"速度流"并与从光流分解得到的运动流对齐，作为数据拟合项加速物理场收敛。整体上，PIDG 实现了物理一致的单目动态重建——不依赖多视角输入，仅凭单目视频即可恢复具有物理意义的粒子运动。

\paragraph{PG-3DGS：物理功能驱动的场景生成}

代表了一条更激进的路线，PG-3DGS（2026.05）~\cite{2605_11266}不再将物理仿真视为 3DGS 的后处理步骤，而是将可微分物理模拟直接嵌入 3DGS 的优化过程中，使生成的场景不仅视觉准确，还满足物理功能目标（如茶壶能倒水、机翼能产生升力）。与 PhysGaussian 的"先重建再仿真"范式不同，PG-3DGS 将物理仿真纳入重建目标函数本身。优化目标 $J(\theta; \mathcal{I}, \mathcal{P}) = L_{\text{vis}}(\theta; \mathcal{I}) + L_{\text{phys}}(\theta; \mathcal{P}) + L_{\text{reg}}(\theta)$ 在视觉损失之外增加物理损失 $L_{\text{phys}} = \mathcal{P}(\Phi(\theta))$，其中 $\Phi(\theta)$ 是从高斯几何到物理仿真网格的转换，$\mathcal{P}$ 是任务相关的物理功能评价（如倒出的液体体积、产生的升力）。

直觉上，视觉损失只能约束外观——两个外观相同的茶壶，一个壶嘴连通内部可以倒水，另一个壶嘴只是外部凸起无法出水，仅靠图像无法区分。

\textbf{动机。}现有 3DGS 重建方法仅优化视觉保真度，无法保证生成物体的物理功能性（如容器是否可盛水、机翼是否产生升力），因此需要将物理仿真约束嵌入重建目标函数中。PG-3DGS 首次在 3DGS 优化中引入这类物理约束，通过可微分流体仿真将物理梯度反传至高斯参数，使优化过程自动调整几何以满足功能需求。生成的场景可直接用于物理交互，无需人工后处理。

紧耦合路线将高斯基元直接作为物理粒子，优势是无需格式转换，但限制为 MPM 一种物理引擎。许多实际场景需要更灵活的物理仿真器——例如刚体碰撞、有限元分析或商业游戏引擎。为此，松耦合路线通过引入代理结构（proxy structure）桥接高斯基元与任意物理引擎。

\paragraph{GASP：代理点结构解放任意黑盒引擎}

与紧耦合路线将高斯基元直接作为 MPM 粒子不同，GASP（2024.09）~\cite{2409_05819}修改了高斯基元与物理引擎的对接方式——用代理点结构替代高斯基元直接参与仿真，从而将物理仿真从 MPM 单一引擎解放为任意黑盒物理引擎。紧耦合路线的局限在于：PhysGaussian 等方法需要修改 MPM 内核以适应高斯参数（如用一阶 Taylor 近似传播变形梯度），因此无法接入 Blender、Genesis 等通用仿真软件。GASP 的核心思路是：物理引擎只关心点的运动，那就只把点的位置交给物理引擎，再把运动结果映射回高斯参数。引擎是黑盒，GASP 不需要知道引擎内部如何计算。

\paragraph*{代理结构：扁平高斯→三角面片→物理顶点}

GASP 采用 GaMeS 参数化（Gaussian-Mesh Splatting，一种扁平高斯表示，即令高斯在某个轴上缩放为 0），将每个 3D 高斯表示为一个三角面片。直觉上，一个扁平高斯就像一张"扑克牌"，由中心点 $\mu$ 和两个沿主轴的边长 $s_2, s_3$ 确定，可以自然地转化为三角形的三个顶点：

\begin{equation}
V = [v_1, v_2, v_3], \quad v_1 = \mu, \quad v_2 = \mu + s_2 \, r_2, \quad v_3 = \mu + s_3 \, r_3,
\end{equation}

其中 $r_2, r_3$ 是旋转矩阵 $R$ 的两个列向量。这三个顶点就是 GASP 交给物理引擎的代理点。反之，物理引擎输出变形后的顶点 $\hat{V}$ 后，GASP 通过逆映射 $\phi^{-1}$ 从 $\hat{V}$ 恢复高斯参数——中心取 $v_1$，方向向量由 Gram-Schmidt 正交化计算，缩放取边长。

\paragraph*{黑盒集成与高斯控制}

由于物理引擎只操作离散点，GASP 能与 Genesis、Blender、Taichi Elements 等任意黑盒引擎无缝集成。引擎输出每步的顶点位置变化 $\hat{V}_t = \{(V_i, t)\}$，GASP 将其映射回高斯参数 $\{(\phi^{-1}((V_i, t)), o_i, c_i)\}_{i=1}^P$。然而，当物理引擎独立处理每个顶点时，某些点可能出现速度波动，导致三角面片急剧膨胀，进而产生巨大的高斯基元伪影。为此，GASP 提出高斯控制规则（Gaussian Control Rules）：当变形后顶点间距超过原始间距的 $\lambda$ 倍时（$\lambda$ 为超参数），自动将缩放值钳制为原始值乘以 $\lambda$，即
\begin{equation}
\| \hat{v}_1(t) - \hat{v}_i(t) \| > \lambda \, \| v_1 - v_i \|, \quad i = 2, 3 \quad \Rightarrow \quad s_i \leftarrow s_i \cdot \lambda.
\end{equation}

此外，GASP 引入 Core-Gaussians 层次化策略：通过 BIRCH（Balanced Iterative Reducing and Clustering using Hierarchies，一种增量式层次聚类算法）聚类将高斯分为 Core-Gaussians（每簇的代表）和 Sub-Gaussians（簇内成员），物理仿真仅在 Core-Gaussians 上执行，然后将运动模式传递给 Sub-Gaussians，显著减少参与仿真的点数。

\paragraph{MaGS：网格吸附高斯的游走式代理}

将 GASP 的固定映射进一步放松，MaGS（2024.06）~\cite{2406_01593}提出网格吸附高斯表示（Mesh-adsorbed Gaussian Splatting），允许 3D 高斯在网格附近游走而非钉死在顶点上，形成互相吸附的网格-高斯混合表示。这一设计的动机在于：GASP 的固定映射在大幅变形时可能产生高斯扭曲或撕裂，而游走式吸附让高斯在保持渲染质量的同时跟随网格运动。网格提供物理仿真的结构化拓扑，高斯提供高质量渲染的灵活性。配合 RMD-Net（残差运动解码器）、RGD-Net（残差几何解码器）和 MPE-Net（运动先验编码器）三个专用网络，MaGS 能从视频中学习运动先验并泛化到新的变形。

\paragraph{PhysTalk：凸包代理与语言驱动仿真}

从 MaGS 的网格吸附策略继续往前，PhysTalk（2025.12）~\cite{2512_24986}进一步简化了代理表示：GASP 使用三角面片作为代理、MaGS 使用吸附网格作为代理，而 PhysTalk 用更轻量的凸包替代两者，并将控制入口从物理引擎参数提升到自然语言。具体而言，PhysTalk 利用 LLM 将自然语言翻译为物理引擎的可执行代码，通过轻量级凸包代理将物理模拟结果传递给 3DGS。整个过程无需训练：LLM 输出代码驱动物理引擎，凸包代理映射物理状态到高斯位置，渲染管线输出最终画面。这是首个将 3DGS 直接与物理模拟器耦合（无需网格提取）的语言驱动框架。

上述方法假设物理属性（质量、弹性、摩擦系数等）已知或手动设定。然而，从视觉观察中自动推断物理属性是实现"全自动物理仿真"的关键一步。物理属性感知路线利用多模态大语言模型（MLLM，Multimodal Large Language Model，包括视觉语言模型 VLM）的先验知识或可微分物理引擎的梯度反传来自动估计物理参数。

\paragraph{Sim Anything：MLLM 零样本属性推断与自适应采样}

Sim Anything（2024.11）~\cite{2411_12789}修改了3DGS物理仿真中的属性设定环节——用MLLM零样本视觉推理替代手动指定或视频扩散模型迭代优化物理参数，首次实现从静态场景到可交互仿真场景的全自动转换。

前述方法面临一个共同瓶颈：PhysGaussian等紧耦合方法需要手动设定材料参数，PhysDreamer等视频扩散方法（通过生成视频并观察振动模式来估计刚度）虽能自动估计但计算昂贵、可控性差，且仅适用于非刚性小物体。人类仅凭视觉就能推断物体的物理属性（金属重而硬、橡胶轻而弹），Sim Anything 正是利用多模态大语言模型（MLLM）模拟这一推理过程。

其流程分为三个阶段。第一阶段，MLLM-P3（MLLM-based Physical Property Perception）对场景中的物体进行物理属性的均值预测：先通过 BLIP 生成物体描述文本，再由 GPT-4V 从图像和文本中推理出 $K$ 个候选材料，最后通过 CLIP 相似度匹配选择最可信材料，输出均值物理属性向量 $\mathbf{M} = (\rho, E, \nu)$（密度、杨氏模量、泊松比）。

第二阶段，MPDP（Material Property Distribution Prediction）将属性预测从逐点回归重构为概率分布估计。直觉上，即便同一材料的不同区域也存在物理属性的固有差异，而 MLLM 只能给出全局均值。MPDP 的核心做法是训练一个小型网络，以物体 3D 点云和 MLLM 预测的均值为输入，预测每个高斯粒子的物理属性分布，再与全局均值逐元素相乘得到每个点的最终属性值。MPDP 训练一个网络 $D_\theta$，以物体的 3D 点云位置和 MLLM 预测的均值为输入，预测每个高斯粒子的物理属性分布 $P = D_\theta(\mathbf{X})$。再将分布 $P$ 与全局均值逐元素相乘，得到每个点的最终属性值。这种"均值缩放分布"的设计使 MPDP 的计算开销仅为 Physics3D 的约 $2\%$。

第三阶段，PGAS（Physical-Geometric Adaptive Sampling，物理-几何自适应采样）解决仿真效率问题。3DGS场景通常包含百万级高斯，直接全部仿真计算量巨大。PGAS 的直觉是：软物体和形状复杂的区域需要更多驱动粒子才能精确模拟形变，而硬物体和平坦区域则不需要。PGAS 自适应调整采样半径 $r$：

\begin{equation}
r_{\text{adapt}} = \text{clip}\!\left(r_0 \cdot k \cdot \frac{E}{\kappa},\; V_{\min},\; V_{\max}\right),
\end{equation}

其中 $E$ 为杨氏模量，$\kappa$ 为物体表面曲率（由高斯位置协方差矩阵的特征值计算），$k=10$，$V_{\min}=1$，$V_{\max}=10$。软材料（$E$ 小）和高曲率区域（$\kappa$ 大）获得更小的采样半径，即更密集的驱动粒子。采样后的粒子通过 MLS-MPM（Material Point Method 的改进版本，采用移动最小二乘近似提升精度）驱动物理仿真。最终整个流程在单 GPU 上约 2 分钟内完成——将静态 3DGS 场景转化为可交互的物理仿真场景。

\paragraph{PUGS：多维度属性扩展}

沿同一方向，PUGS（2025.02，Physics Understanding via Gaussian Splatting）~\cite{2502_12231}聚焦于 Sim Anything 未充分探索的属性维度扩展——Sim Anything 仅预测密度、杨氏模量和泊松比三个力学属性，而 PUGS 利用 3DGS 重建物体并结合 VLM 零样本预测多种物理属性（质量、摩擦系数、硬度），引入几何感知正则化损失和区域感知特征对比损失提升形状质量与局部亲和性，在 ABO-500 质量预测基准上达到新 SOTA。

\paragraph{PhysGS：贝叶斯不确定性量化}

针对 Sim Anything 的 MLLM 点估计缺乏不确定性量化这一局限，PhysGS（2025.11）~\cite{2511_18570}将物理属性估计建模为 3DGS 上的贝叶斯推理问题：使用 Dirichlet-Categorical 模型融合多视角材料标签，Normal-Inverse-Gamma 先验估计连续物理属性并建模认知不确定性和偶然不确定性，质量估计精度提升 22.8\%，硬度误差降低 61.2\%。

\paragraph{PIXIE：前馈通用属性预测网络}

与 Sim Anything 依赖 MLLM 零样本推理（测试时优化、单场景不泛化）的范式不同，PIXIE（2025.08）~\cite{2508_17437}更进一步，提出首个通过监督学习训练通用神经网络预测 3D 场景物理属性的前馈框架——将属性预测从"每次推理"提升为"一次训练、处处适用"。其核心洞察是利用 CLIP（Contrastive Language-Image Pre-training，对比语言-图像预训练）等预训练视觉特征将外观与物理特性关联——金属看起来光滑且反光，橡胶看起来粗糙且哑光。PIXIE 构建了 PIXIEVERSE 数据集（最大的 3D 资产与物理标注配对数据集），比测试时优化方法快数个数量级且效果提升 1.46--4.39 倍。

\paragraph{FastPhysGS：加速端到端物理仿真}

在 PIXIE 将属性预测从优化范式转向前馈范式之后，后续工作沿加速和输入简化两个方向展开。FastPhysGS（2026.02）~\cite{2602_01723}针对 Sim Anything 约 2 分钟的端到端流程，聚焦于加速物理仿真流程：提出实例感知粒子填充（IPF）使用蒙特卡洛重要性采样高效填充内部粒子，双向图解耦优化（BGDO）通过自适应策略快速优化 VLM 预测的材料参数，在 1 分钟内完成高保真物理仿真。

\paragraph{ReconPhys：单目视频联合推断属性与几何}

与 Sim Anything 需要预建 3DGS 场景再推断属性的两步流程不同，ReconPhys（2026.04）~\cite{2604_07882}提出首个前馈框架从单目视频联合推断物理属性和 3DGS，采用弹簧-质点系统绑定 3D 高斯，推理速度 $<$1 秒，未来预测 PSNR 达 21.64（vs.\ 优化基线的 13.27）。

\paragraph{PIG：多材料物理交互}

上述方法均假设场景中只有单一材料，PIG（2025.06，Physics-aware Interaction of Gaussians）~\cite{2506_07657}将问题推进到 Sim Anything 未触及的多材料物理交互：通过 2D 像素到 3D 高斯的快速精确映射实现物体级分割，为不同物体分配独特物理属性进行多材料耦合交互，并提出基于约束尺度的变形梯度方法消除渲染伪影。

\paragraph{Mirage2Matter：从单物体仿真到完整世界模型}

将视野从单物体仿真扩展到完整世界模型，Mirage2Matter（2026.02）~\cite{2602_00096}与 Sim Anything 停留在静态场景的物理化不同，将 3DGS 重建与生成模型结合构建物理真实的仿真世界模型，利用标定靶实现重建场景与真实世界的精确尺度对齐，VLA（Vision-Language-Action，视觉-语言-动作模型）模型在该仿真数据上训练后实现零样本性能。其核心做法是：先用标定靶将 3DGS 重建场景与真实世界精确尺度对齐，再利用生成模型补全缺失区域，最终构建可供机器人策略训练的仿真世界模型。

\subsection{流体与粒子动力学}

流体是物理仿真中最具挑战性的材料类型之一——不可压缩性、涡旋结构和自由表面等特征对空间表示提出特殊要求。3DGS 的离散高斯基元天然适合作为拉格朗日粒子跟踪流体运动，但其各向异性椭球形状与流体的连续速度场之间存在表示鸿沟。本节的方法沿两条路线展开：一条从流体模拟器出发，将高斯基元用作流体粒子；另一条从重建出发，从视频观测中恢复流体动力学。下面先介绍高斯空间表示与无网格流体求解，再讨论从视频恢复流体动力学的方法。

\subsubsection{高斯空间表示与无网格流体求解}

本路线介绍利用3DGS的高斯空间表示构建无网格流体求解器以及从视频观测中恢复流体动力学的方法。

\paragraph{Gaussian Fluids：无网格流体求解器}

Gaussian Fluids（2024.05）~\cite{2405_18133}首次将 3DGS 的高斯空间表示引入流体模拟领域——用连续高斯核加权和替代传统的欧拉网格或拉格朗日粒子离散化，构建无网格（grid-free）流体求解器。这也是本路线的代表性工作。

\textbf{Baseline 问题。}传统流体求解器依赖两种空间表示：欧拉方法在固定网格上存储速度值，受限于网格分辨率且容易引入数值黏性（插值导致非物理耗散）；拉格朗日方法（如 SPH，Smoothed Particle Hydrodynamics，光滑粒子流体动力学）用粒子跟踪流体运动，但精度有限，难以捕捉精细涡旋结构。

\textbf{动机。}传统网格/粒子混合方法在两种离散化之间传递数据时引入数值误差，而隐式神经表示虽提供连续场但优化收敛慢，因此需要一种兼具连续可微与高效的空间表示。混合方法虽试图结合两者优势，但在两种离散化之间传递数据时引入数值误差，且需维护双重表示的额外开销。隐式神经表示（INR，Implicit Neural Representation）虽提供连续表示，但优化收敛慢、难以施加硬物理约束。

\textbf{核心表示：高斯空间表示（GSR，Gaussian Spatial Representation）。}Gaussian Fluids 的关键洞察是：将流速场视为一堆"柔软的高斯斑点"的叠加——每个高斯核覆盖一块局部区域，携带一个速度向量，所有核的贡献加起来构成完整的速度场。直觉上，这就像用一组可变形、可移动的"速度画笔"来描绘流场，而非用固定网格点。具体地，连续流速场建模为多个高斯函数的加权和：

\begin{equation}
\tilde{\mathbf{v}}(\mathbf{x}) = \sum_{i=1}^{N} \mathbf{v}_i \, \hat{G}_i(\mathbf{x}),
\end{equation}

其中 $\hat{G}_i$ 为第 $i$ 个截断高斯核（在远离中心处截断为零以加速计算），$\mathbf{v}_i \in \mathbb{R}^m$ 为该核携带的权重向量。每个高斯核自身有参数 $\Theta = \{\mathbf{v}_i, \boldsymbol{\mu}_i, \theta_i/\mathbf{r}_i, \mathbf{s}_i^{-1}\}$（权重、位置、旋转、逆尺度），通过优化来拟合和演化流场。也就是说，流场在任意空间点 $\mathbf{x}$ 处的速度，等于所有"覆盖"该点的高斯核的贡献之和。

\textbf{求解流程。}Gaussian Fluids 遵循经典的算子分裂框架，每个时间步分三阶段：(1) \emph{重播种}（Reseeding）——将过度拉伸的高斯粒子分裂（最大尺度超过最小尺度的 $r_\mathrm{aniso}=1.5$ 倍时分裂），补充自适应分辨率；(2) \emph{平流初始猜测}——将高斯粒子视为拉格朗日粒子，沿速度场用四阶 Runge-Kutta 平流至新位置，作为优化的初始猜测；(3) \emph{物理引导优化}——以涡度损失 $\mathcal{L}_\mathrm{vor}$（匹配平流后的涡旋场）和散度损失 $\mathcal{L}_\mathrm{div}$（施加不可压缩约束 $\nabla \cdot \mathbf{v} = 0$）为核心，辅以各向异性正则 $\mathcal{L}_\mathrm{aniso}$（限制高斯长宽比 $\leq 1.5$）、体积正则 $\mathcal{L}_\mathrm{vol}$（保持粒子体积均匀）和位置惩罚 $\mathcal{L}_\mathrm{pos}$（防止粒子聚集），联合优化所有参数。一个关键技术细节是梯度投影：当涡度损失和散度损失的梯度方向矛盾时，将两者投影到彼此正交的方向上，避免梯度冲突导致的振荡伪影。

\textbf{GSR 的三重优势。}其一，内存高效——不需要密集网格，高斯自适应分布到需要精细分辨率的区域。其二，连续可微——速度场在任意位置都有定义且可直接解析求导，无需自动微分，无网格插值误差。其三，空间自适应——高斯可以稀疏分布在低涡旋区域、密集分布在需要精细解析的高涡旋区域。相比隐式神经表示，GSR 在稳定性、精度和跨场景泛化性上均更优。该方法在 Taylor 涡旋、Kármán 涡街、3D 涡环碰撞等多种 2D/3D 流体场景中验证了精细涡旋结构的保持能力，且无需逐场景调参。

\paragraph{LagrangianSplats：无散度核约束的流体重建}

与 Gaussian Fluids 从\emph{模拟器设计}出发、将高斯作为流场表示的主动组件不同，LagrangianSplats（2026.05）~\cite{2605_09299}从\emph{流体重建}的角度提出结构化约束框架：用无散度核（Divergence-Free Kernel, DFK）参数化速度场，\emph{在构造上}保证速度场严格无散度（不可压缩），而非通过惩罚项软约束。直觉上，DFK 将不可压缩性"嵌入"表示本身——无论参数如何变化，速度场始终满足 $\nabla \cdot \mathbf{v} = 0$。

高斯基元作为被动示踪物被速度场驱动平流，位置由拉格朗日平流算子确定 $\mathbf{x}_i^t = \mathcal{A}_{0 \to t}(\mathbf{x}_i^0, \mathbf{v})$。为了处理长序列优化中的梯度消失问题，LagrangianSplats 引入滑动窗口策略：将长视频分割为重叠的时间窗口，在每个窗口内联合优化速度场和高斯属性，窗口之间通过连续的平流轨迹保证物理一致性。

\paragraph{FluidGaussian：仿真不确定性引导的智能重建}

与 Gaussian Fluids 和 LagrangianSplats 均以\emph{仿真为核心目标}不同，FluidGaussian（2026.03）~\cite{2603_21356}转向\emph{重建}视角，提出了一个全新的思路：将流体仿真的\emph{不确定性}引入 3DGS 重建。传统重建方法追求视觉质量，但重建结果在物理上可能不可行（如流体内部存在非零散度）。FluidGaussian 定义基于流体仿真的不确定性度量，将其与主动学习结合，优先选择同时提升视觉和物理保真度的视角。在 NeRF Synthetic、Mip-NeRF 360 和 DrivAerNet++ 上实现 PSNR 提升 8.6\%、流体仿真速度散度降低 62.3\%，开创了"功能智能 3D 重建"方向。

\paragraph{FLUIDSPLAT：广义物理场的高斯基元重建}

与前述方法均围绕\emph{流体本身}展开不同，FLUIDSPLAT（2026.05）~\cite{2605_18866}将 3DGS 的空间表示思想迁移至广义物理场重建（非流体仿真）：从稀疏表面传感器重建连续物理流场（速度场、压力场）。提出用 $K$ 个各向异性高斯基元构成 partition-of-unity 支架作为空间显式中间表示，配合状态条件残差解码器提供支架无法捕获的细节。理论上证明了高斯基元数量的最优缩放律 $K^* \sim (N/\sigma^2)^{d/(2s+d)}$，在四个 2D/3D 基准上实现 11--28\% 误差降低。

上述方法从模拟器或重建的角度"正向"处理流体问题——设计求解器或从视频恢复速度场。另一个方向则更为实用：从单目视频直接提取可重仿真的 3D 流体资产。

\paragraph{从视频恢复流体动力学}

3D Dynamic Fluid Assets（2025.03）~\cite{2503_00868}修改了流体重建的表示与优化环节——用\emph{可微分物理仿真}替代纯视觉重建，使输出不仅是视觉资产，而是"仿真就绪"的物理流体资产。传统单目 3D 重建方法（如 NeRF 或标准 3DGS）仅恢复外观和几何，无法提取流体的速度场、密度、黏度等物理属性，因此重建结果不可重仿真，也无法编辑交互。该方法正是为了解决这一问题：从单目视频中提取几何、速度场和仿真参数，输出可直接导入物理引擎的流体资产。

整个流程分为三步，核心思路是\textbf{先用视觉方法获取几何和表面速度，再用物理方法优化内部体积速度和仿真参数}：

\textbf{第一步：生成式 3DGS 获取与预处理流体几何。}利用图像条件的 3DGS 生成模型（如 TriplaneGaussian）从单帧图像生成初始 3DGS，再通过单视图优化提升分辨率。优化损失函数包含五项：图像重建损失 $L_\text{img}$、各向异性损失 $L_\text{aniso}$（抑制过度拉伸）、体积损失 $L_\text{vol}$、尺度损失 $L_\text{scl}$（鼓励更细粒度的高斯基元）和亮度损失 $L_\text{lumi}$（惩罚过饱和）。优化后进行去噪和体素化填充：先剪枝低透明度或过度拉伸的高斯基元，再以网格采样填充流体内部空腔（通过射线-壳体交点奇偶性判断内部点），确保流体体积完整。直觉上，这一步将生成模型产生的稀疏"壳"变为密实的"体"。

\textbf{第二步：从光流到体积速度场。}利用光流估计 2D 屏幕速度 $V_\text{2D}$，结合 3DGS 的深度信息反投影到 3D 表面速度 $V_\text{surf}$。然而，光流在流体局部纹理均匀的区域（如平静水面）往往检测失败，输出接近零速度。为此，该方法提出\emph{主流校正}（Mainstream Correction）策略：(1) 利用几何约束（如河岸走向）确定主流方向，从邻近有效光流区域插值补偿缺失区域；(2) 基于不可压缩约束推导 2D 散度方程 $\nabla_\text{2D} \cdot V_\text{2D} = -(\frac{1}{z}\frac{\partial v_z}{\partial z} + \frac{2v_z}{z}) - v_z \frac{\partial}{\partial u} - v_z \frac{\partial}{\partial v}$，通过投影法保证校正后速度满足物理约束。该方程仅依赖深度方向速度分量 $v_z$（由深度信息确定），在局部平行于相机的水面区域 2D 散度趋于零，物理直觉是：平坦水面的 2D 投影天然接近不可压缩。

获得表面速度后，还需计算整个体积的内部速度场 $V_\text{vol}$。流体底部（如河床）速度未知，该方法利用 CFD 中的壁面函数估算边界层阻尼，采用层流外流的三阶多项式近似：$\frac{|v_\text{base}|}{V_\infty} = \frac{3}{2}\left(\frac{y}{\delta}\right) - \frac{1}{2}\left(\frac{y}{\delta}\right)^3$，其中 $V_\infty$ 为对应表面速度幅值，$\delta$ 为边界层厚度（由流体类型动态决定：河流取较大厚度，烟雾取亚网格尺度），$y$ 为到壁面的距离。该公式直觉含义是：越靠近固体壁面，流体速度越慢，呈现平滑的阻尼过渡。最后，以封闭表面速度为边界条件，通过约束投影法求解无散度 3D 速度场。

\textbf{第三步：可微分仿真参数优化。}将粒子速度场转换到欧拉网格上，利用 Affine Particle-in-Cell（APIC，一种混合拉格朗日-欧拉数值方法，通过额外计算局部仿射速度场提升精度和稳定性）建立可微分仿真管线。两个关键修改使仿真可微分：(1) 用 PeRCNN（物理编码循环卷积网络）替代传统迭代压力求解器（如 Jacobi 迭代），学习卷积核近似空间微分算子（梯度、散度、拉普拉斯），外层循环结构模拟迭代过程，在保持精度的同时实现端到端可微分；(2) 对流项用解析梯度替代不可微的回溯网格索引操作。优化目标为 11 个物理参数：入口/出口体积速度 $v_\text{in}$/$v_\text{out}$、入口速度波动 $\sigma_\text{in}$、密度 $\rho$、黏度 $\mu$、边界弹跳 $b$ 与阻尼 $d$、重力 $g$ 和时间步长 $dt$。损失函数结合速度 $L_2$ 距离和方向一致性：$L = \mathbf{1} \cdot \left[(1 - \frac{v_\text{sim}}{|v_\text{sim}|} \cdot \frac{v_\text{gt}}{|v_\text{gt}|}) + L_2(v_\text{sim}, v_\text{gt})\right]$，按网格类型（表面/占据/空）加掩码。

输出为与源视频物理动态匹配的"仿真就绪"流体资产：3DGS 几何 $+$ 优化后的物理参数。该资产支持几何编辑（修改河道形状后重新仿真）、运动时长扩展（仿真可运行比视频更长时间）、以及与外部场景的物理交互。APIC 仿真中，仿射矩阵 $C$ 近似速度梯度，通过变形梯度 $F_p^{n+1} = (I + dt \cdot C_p^{n+1}) F_p^n$ 更新高斯协方差 $A_p = F_p A_p^0 F_p^\top$，使高斯基元的形状随流体变形自然演化。

\subsection{可变形物体与布料仿真}

可变形物体（绳子、布料、毛绒玩具等）的动力学建模是物理仿真的重要子领域。与刚体和流体不同，可变形物体既有整体运动（如布料飘动），又有局部形变（如布料褶皱），且不同部位可能呈现截然不同的材料行为（弹性、黏弹性、塑性等）。3DGS 的挑战在于：高斯基元缺乏拓扑连接关系，无法直接表达可变形物体的内部结构。本节从弹簧-质点图结构化物理骨架出发，逐步讨论该领域的代表性方法。

\subsubsection{弹簧-质点图结构化物理骨架}

本路线介绍用弹簧-质点图等结构化物理骨架为离散高斯点注入拓扑连接关系的方法，使可变形物体的动力学建模具备物理约束和外推能力。

\paragraph{PhysTwin：弹簧-质点图的数字孪生}

PhysTwin（2025.03）~\cite{2503_17973}修改了可变形物体的动力学表示环节——用弹簧-质点图结构替代标准 3DGS 的无拓扑高斯基元集合，为离散高斯点注入结构化的物理连接关系。标准 3DGS 的动态变体（如 Deformable 3DGS、Dynamic 3DGS）通过学习形变场或逐帧优化高斯位置来拟合运动，但缺少物理约束使其无法外推到未见交互场景（如改变抓取方式或施加新的外力）。PhysTwin 的设计动机正是针对这一问题：以弹簧-质点模型作为物理骨架，每个质点通过弹簧与相邻质点连接，弹簧的弹力提供恢复力，阻尼力耗散能量，从而在拓扑层面建立可变形物体的内部结构。

PhysTwin 的核心表示是一个弹簧-质点图 $G=(V,E)$，其中 $V$ 为质点集、$E$ 为弹簧集。质点 $i$ 受到相邻弹簧力和阻尼力的合力驱动：

\begin{equation}
\mathbf{F}_i = \sum_{(i,j)\in E} \left[ k_{ij}(\|\mathbf{x}_j - \mathbf{x}_i\| - l_{ij})\frac{\mathbf{x}_j - \mathbf{x}_i}{\|\mathbf{x}_j - \mathbf{x}_i\|} - \gamma(\mathbf{v}_i - \mathbf{v}_j) \right] + \mathbf{F}_i^{\text{ext}}
\end{equation}

其中 $k_{ij}$ 为弹簧刚度，$l_{ij}$ 为弹簧自然长度，$\gamma$ 为阻尼系数，$\mathbf{F}_i^{\text{ext}}$ 包含重力、碰撞等外力。系统通过显式欧拉积分更新状态 $\mathbf{X}_{t+1} = f_{\boldsymbol{\theta}, G_0}(\mathbf{X}_t, \mathbf{a}_t)$，其中 $\boldsymbol{\theta}$ 编码所有物理参数（刚度、阻尼、碰撞参数），$G_0$ 为初始几何与拓扑，$\mathbf{a}_t$ 为控制点动作。直观而言，弹簧力将质点拉回平衡位置（类比橡皮筋），阻尼力防止永久振荡（类比减震器），二者共同约束物体的弹性变形行为。

PhysTwin 的核心流程分两个阶段。第一阶段聚焦物理与几何优化：先用生成式形状先验（TRELLIS）从单张 RGB 图像恢复完整几何——因为稀疏视角下深度观测仅覆盖部分表面，直接从点云恢复完整拓扑不可行——获得初始网格后，通过配准模块（基于 SuperGlue 的 2D 匹配估计旋转、PnP 精修位姿、as-rigid-as-possible 非刚性形变）对齐到真实观测。随后以稀疏到密集的分层优化策略联合优化拓扑和物理参数：先以零阶采样优化（CMA-ES，Covariance Matrix Adaptation Evolution Strategy，协方差矩阵自适应进化策略）估计不可微的拓扑参数（连接半径、最大邻居数）和均匀刚度，再用一阶梯度下降通过自建可微弹簧-质点仿真器精修稠密刚度（单个弹簧独立刚度，参数空间达数万维）。

第二阶段进行外观优化：以线性混合蒙皮（LBS，即根据邻近质点的运动加权插值高斯基元的位置和形状）将高斯绑定到质点网格，使高斯随质点运动自动变形。外观通过 $L_1$ 损失加 D-SSIM 项在首帧优化，高斯被约束为各向同性以避免形变时的尖刺伪影。

PhysTwin 的关键设计是引入控制点——通过弹簧将控制点连接到物体点（参数为连接半径和最大邻居数），使系统支持交互式仿真和基于模型的机器人运动规划。在 22 个真实可变形物体场景（绳索、毛绒玩具、布料、快递盒）上的实验表明，PhysTwin 在重建重仿真、未来预测和未见交互泛化三个任务上均优于 Spring-Gaus 和 GS-Dynamics 基线，且支持实时交互仿真（基于 Warp 加速）。

\paragraph{NeuMA：物理先验与学习修正的分解策略}

与 PhysTwin 以弹簧-质点图作为显式物理骨架不同，NeuMA（2024.10）~\cite{2410_08257}转而从视觉观察中学习物体内在动力学，提出了一种优雅的分解策略：将材质模型分解为物理先验和可学习修正项 $\sigma = \sigma_{\text{prior}} + \sigma_{\text{correct}}$，其中 $\sigma_{\text{prior}}$ 来自经典本构模型（如 Neo-Hookean 弹性），$\sigma_{\text{correct}}$ 由神经网络预测。PhysTwin 的弹簧参数通过零阶采样和一阶梯度联合优化确定，而 NeuMA 将物理知识嵌入为结构化先验（而非可优化参数），使学习过程聚焦于修正项。这一分解兼顾了物理先验的泛化性（对已知材料行为建模准确）和学习修正的灵活性（捕获偏离理想模型的真实行为），配合 Particle-GS 可微分渲染器实现端到端优化。

\paragraph{3DGSim：纯神经力场的端到端物理模拟}

与 NeuMA 将物理先验嵌入为结构化约束的``先验+修正''范式不同，3DGSim（2025.03）~\cite{2503_24009}进一步放弃了显式物理模型的依赖，提出首个完全端到端可微的 3D 物理模拟框架，从多视角 RGB 视频联合训练逆渲染和动力学预测，无需深度信息或粒子追踪等特权信息。PhysTwin 和 NeuMA 均保留了显式物理结构（弹簧-质点图或本构模型），而 3DGSim 用纯神经力场替代物理方程：学习全局变换力场 $F^{\text{global}}$（捕获物体间接触、滑动和重力）和局部应力场 $F^{\text{local}}$（捕获软体变形），通过二阶积分器恢复物理轨迹。该框架能统一建模从刚性到弹性、布料等多种物理行为。

\paragraph{Particle-Grid Neural Dynamics：拉格朗日-欧拉混合动力学}

与 3DGSim 用纯神经力场建模动力学不同，Particle-Grid Neural Dynamics（2025.06）~\cite{2506_15680}回归混合表示的思想，提出拉格朗日-欧拉混合表示学习可变形物体动力学：粒子表示携带物体状态（拉格朗日视角，跟踪个体运动），均匀网格注入空间连续性先验（欧拉视角，提供全局空间信息）。其关键组件包括全局点云编码器（PointNet）捕获整体形状信息以克服 GNN 的连通性限制，网格速度编辑（GVE）机制建模物体-地面碰撞和机器人夹持，以及 B 样条插值实现网格到粒子的速度传递。结合 3DGS 的 LBS 绑定实现可变形物体的数字孪生和动作条件视频生成。

\paragraph{MonoPhysics：单目逆物理估计}

上述方法均依赖多视角观测，而 MonoPhysics（2026.05）~\cite{2605_30320}将可变形物体动力学建模推进到更具挑战性的单目设置，实现了首个单目逆物理估计框架，通过三个创新的视觉-物理桥接机制将可微分 MPM 模拟与 3DGS 耦合：(1) 尺度重参数化——在相机空间维护粒子并通过单一可学习标量 $s$ 控制全局尺度，使图像损失通过可微分仿真聚合一致的尺度信号；(2) 物理感知几何精修——使每个粒子的物理体积 $V_i$ 成为可优化参数并通过分区-of-unity 约束自适应更新，配合物理重要性驱动的粒子管理；(3) 可微分位置图——通过重参数化 $\hat{x}_i = x_i^{2D} + (\Sigma_i^{2D})^{1/2} z$ 使像素位置成为高斯参数的可微函数，通过 Sinkhorn 散度的最优传输提供全局方向梯度。在单目设置下达到多目方法的性能水平。

除上述代表性方法外，可变形物体动力学建模还在运动表示设计、物理基元构造和布料动力学三个方向上持续拓展。下面分别介绍这三个方向的进展。

在\emph{运动表示}方向，核心挑战是如何在保持计算效率的同时精确捕捉复杂物体的全局和局部运动。Dynamic 3D Gaussian Tracking（2024.10）~\cite{2410_18912}结合动态 3DGS 重建与图神经网络（GNN，Graph Neural Network）从真实多视角视频中学习可变形物体动力学。

\textbf{动机。}密集高斯基元数量庞大，直接在其上训练动力学模型计算代价极高；该方法通过稀疏控制粒子作为中间表示，使动力学模型既学习全局交互模式又保持实时推理能力。其关键设计是将密集 3D 高斯重建结果下采样为稀疏控制粒子（control particles），在控制粒子上运行 GNN 动力学模型预测运动，再通过插值将运动传播回密集高斯。这种"稀疏驱动、稠密插值"的策略使动力学模型既能学习到全局的物体间交互模式，又保持了可部署于机器人操控规划的实时推理能力。HEIR（2025.10）~\cite{2510_26786}则从运动分解的视角出发，提出一种通用的层次运动建模方法：通过可微图学习自动发现运动的层次结构，将每个运动单元视为图的顶点，有向边编码学习到的父子依赖关系，从而将全局绝对运动显式分解为父级继承模式（parent-inherited pattern，即子节点从父节点继承的运动分量）和局部运动残差（local motion residual，即子节点独有的精细运动）。这种数据驱动的层次分解无需手工定义运动基元，可在不同任务间泛化，且分解出的层次结构本身具有可解释性。

在\emph{物理基元}方向，PhysConvex（2026.02）~\cite{2602_18886}跳出高斯/粒子的传统表示范式，提出用物理驱动的 3D 动态凸面基元（dynamic convex primitives）统一视觉渲染和物理模拟。

\textbf{动机。}标准 3DGS 的椭球基元缺乏物理属性，需要外部仿真器驱动运动；PhysConvex 将连续力学直接嵌入凸面基元本身，使视觉渲染和物理仿真在统一表示上完成，避免了表示转换开销。与标准 3DGS 用无物理属性的椭球体表示场景不同，PhysConvex 的凸面基元由连续力学（continuum mechanics）支配，通过边界驱动动态凸面表示（boundary-driven dynamic convex representation）建模变形：顶点和面的动力学演化捕捉空间自适应、非均匀的变形行为以及动态变化的边界。为进一步提升大规模复杂几何和异质材料的仿真效率，PhysConvex 开发了降阶凸面模拟（reduced-order convex simulation）策略——利用神经蒙皮特征模态（neural skinning eigenmodes）作为形状和材料感知的变形基底，在时间变化的参数空间中驱动凸面场的对流（advection）。这一策略在保持物理精度的同时显著降低了计算维度。

在\emph{布料动力学}方向，CloDS（2026.02）~\cite{2602_01844}聚焦布料这一具有大非线性变形和严重自遮挡的特定可变形物体类别，提出在未知物理条件下从纯视觉观测无监督学习布料动力学。

\textbf{动机。}现有方法依赖已知物理属性（如刚度、阻尼）作为监督或输入，但这些参数在真实场景中难以获取；CloDS 采用无监督策略，从视频直接学习布料动力学，无需任何物理标注。为应对接地过程中布料的大幅非线性变形和自遮挡问题，CloDS 引入双位置不透明度调制（dual-position opacity modulation）机制——同时考虑绝对位置和相对位置的网格高斯泼溅，实现 2D 观测与 3D 几何之间的双向映射。这使得系统在无需任何物理标注的条件下也能学习到物理一致的布料运动模式。

\paragraph{布料仿真专用方法}

布料是一类特殊的可变形物体，具有薄壳结构和大变形特征，其状态估计面临无穷维状态空间和严重自遮挡的双重挑战。

\textbf{Baseline 做法与问题。}现有布料状态估计主要依赖两类方法：(1) 基于深度传感器的 3D 表示（如点云、图结构），虽能捕捉几何但缺乏纹理等视觉线索，且仿真到真实的迁移需要耗时的在线微调；(2) 基于 3DGS 的动态跟踪方法（如 MD-Splatting、DeformGS）需要至少 50 个相机的密集观测和逐场景优化，计算开销极大，难以满足实际部署需求。

Cloth-Splatting（2025.01）~\cite{2501_01715}修改了布料状态估计的观测模型环节——用\emph{mesh 约束的 Gaussian Splatting}替代传统的深度传感器或密集多目 3DGS，仅需 3--4 个 RGB 相机即可从图像空间反向传播梯度来修正布料 3D 状态。

\textbf{实现原理。}该方法将布料状态估计建模为贝叶斯滤波问题——类似卡尔曼滤波的``预测-更新''两步循环，其中``预测''由学习到的动力学模型提供先验，``更新''由可微渲染器利用观测修正先验。

预测步骤使用基于图神经网络（GNN）的动作条件动力学模型 $f_\theta$：给定过去 $m$ 帧的布料网格状态 $M_{t-m:t}$ 和机器人动作 $a_t$，预测下一帧状态 $\hat{M}_{t+1} = f_\theta(M_{t-m:t}, a_t)$。网络输出每个顶点的加速度，经前向欧拉积分得到下一帧的速度和位置。

更新步骤的关键创新在于\emph{mesh 约束的 Gaussian Splatting}（$h_{\text{GS}}$）：将每个 3D 高斯绑定到布料 mesh 的一个三角面上，其位置用重心坐标参数化：

\begin{equation}
  \mu = b_1 v^1 + b_2 v^2 + b_3 v^3, \quad b_1 + b_2 + b_3 = 1
\end{equation}

其中 $v^1, v^2, v^3$ 为所属面的三个顶点。直觉上，高斯``粘贴''在 mesh 表面，随布料变形而自动移动——mesh 变形时高斯的重心坐标和外观参数保持不变，仅顶点位置变化驱动高斯重新定位。这使得布料状态空间到图像空间形成完全可微的映射。

为避免直接优化 GNN 导致梯度消失，更新后的状态定义为预测状态加一个学习到的残差偏移：

\begin{equation}
  \tilde{M}_{t+1} = \hat{M}_{t+1} + \delta_{t+1}, \quad \delta_{t+1} = u_\phi(t+1)
\end{equation}

其中残差网络 $u_\phi$ 是一个以归一化时间为输入的 3 层 MLP（宽度 256），输出初始值接近零。通过最小化观测损失 $L_{\text{obs}} = \|Y_{t+1} - h_{\text{GS}}(\tilde{M}_{t+1}, P)\|_2^2$ 优化残差，即让渲染的布料图像与真实 RGB 观测尽可能一致。正则化损失 $L_{\text{reg}} = L_{\text{SSIM}} + L_{\text{iso}} + L_{\text{magn}}$ 分别约束结构相似性、邻域顶点等距（保持布料不可拉伸性）和最小运动（防止数值不稳定）。

\textbf{效果。}相比纯物理基线（仅用 GNN 预测），GS 更新将跟踪误差降低 57\%；相比次优基线 MD-Splatting，收敛速度提升 85\%，且仅需约 4k 个高斯即可重建布料外观。该方法进一步支持闭环折叠操控：通过模型预测控制（MPC）在每次重规划时用 GS 修正布料状态估计，折叠精度接近有真值状态的 Oracle 方法。

SketchRodGS（2026.01）~\cite{2601_02072}从用户交互的角度解决 3DGS 场景中细长几何的提取与动画化问题。其核心思路十分直觉：用户在屏幕上画一条线（草图），系统自动找到对应的 3D 细长结构并将其转化为可仿真的折线网格（polyline mesh）。由于 3DGS 缺乏拓扑连接信息且高斯基元的配置包含大量噪声，直接从点云提取细长结构极为困难。SketchRodGS 的关键创新在于屏幕空间最短路径分析：将用户草图投影到多视角的深度图上，在屏幕空间中通过动态规划高效求解最短路径，从而在不同视角间鲁棒地重建出绳索、布料边缘等细长部件的折线表示。提取的折线可直接接入弹性杆（elastic rod）物理仿真器，使 3DGS 场景中的细长物体能够进行物理驱动的动画——绳索可以摆动、打结，布料边缘可以随风吹动。该方法在多个自然场景（in-the-wild）示例上验证了有效性。

前述方法主要处理单时间步或短序列的物理行为。将物理仿真扩展到完整的时空 4D 场景，则需要解决材料参数的自动推断和跨帧运动一致性两个更根本的问题。下面介绍该方向的代表性方法。

\paragraph{PhysFlow：光流引导的七种材料参数优化}

将物理仿真扩展到 4D（3D + 时间）需要解决参数推断和运动监督两个核心问题。PhysFlow（2024.11）~\cite{2411_14423}修改了材料参数优化的监督信号环节——用光流损失（optical flow loss）替代传统的渲染损失（render loss）和分数蒸馏采样损失（SDS loss），实现从任意输入（单图/多视角/视频）自动估计材料属性并进行 4D 物理仿真。

\textbf{Baseline 做法与问题。}在 PhysFlow 之前，基于视频扩散的材料参数优化主要有两条路线。PhysDreamer 使用渲染损失 $L_{\text{render}} = \sum_t \lambda L_1(\hat{I}_t, I_t) + (1-\lambda)L_{\text{D-SSIM}}(\hat{I}_t, I_t)$，直接比较仿真渲染帧与扩散生成帧的像素差异——但由于视频扩散输出的不稳定性与高噪声，该损失仅能处理微小运动，稍大的运动即导致优化发散。另一条路线（如 Physics3D）采用 SDS 损失 $L_{\text{sds}}$，从预训练扩散模型中蒸馏梯度，但 SDS 需要对整个扩散模型做反向传播，显存开销极大，且在结合复杂文本提示的视频扩散模型时更为严重。此外，上述方法均局限于弹性材料，无法处理流体、沙子等复杂材料。

\textbf{光流引导的参数优化。}PhysFlow 的核心洞察是：优化材料参数并不需要逐像素对比图像，而只需对比\emph{运动场}——光流恰好编码了像素级位移信息，是材料参数对运动的"直接传感器"。具体流程为：首先利用 GPT-4 从输入图像推断材料类型 $M$、密度 $\rho$ 和初始参数 $M_p^0$；然后根据输入类型（多视角/视频/单图）分别选择 PGSR、Grid4D 或 Splatt3R 构建 3D 高斯；接着通过视频扩散模型（根据输入类型可选 SVD、CogVideoX、MotionCtrl 等）生成参考视频，并用 RAFT 计算相邻帧光流 $U(I_t, I_{t+1})$；最终在可微分 MPM（Material Point Method，一种混合拉格朗日-欧拉的粒子仿真方法，将材料离散为粒子追踪变形，在背景网格上求解动量守恒）框架下，最小化光流损失：

\begin{equation}
    L_{\text{flow}} = \sum_t \|U(I_t, I_{t+1}) - \hat{U}(\hat{I}_t, \hat{I}_{t+1})\|^2
\end{equation}

其中 $\hat{U}$ 是仿真渲染帧之间的光流。直观地说，该损失迫使仿真运动与参考视频中的运动对齐，而非逐像素对比外观——这避免了图像级噪声的干扰，同时使梯度信号更直接地反映材料参数对运动的影响。

\textbf{材料类型覆盖。}得益于 MPM 的通用性，PhysFlow 支持 7 种材料类型（弹性、橡皮泥、金属、泡沫、沙子、牛顿流体、非牛顿流体），每种材料由不同参数组合定义：弹性体由杨氏模量 $E$ 和泊松比 $\nu$ 控制；橡皮泥和金属额外引入屈服应力 $\sigma_Y$ 以建模塑性变形；沙子由摩擦角 $\phi_{\text{fric}}$ 表征；流体则由粘度 $\eta$ 和体积模量 $\kappa$ 控制。优化后的参数集成回可微分 MPM 框架，驱动 3D 高斯形变并从新视角渲染，完成 4D 重建。在合成数据集上，PhysFlow 在多数材料类型的相对误差（RE）和端点误差（EPE）上优于 PAC-NeRF，与 GIC 表现相当；在真实数据集上，32 人的人类评估中物理真实感和视觉真实感评分均超过 3 分（5 分制），显著高于所有基线方法。

\paragraph{NGFF：端到端神经力场预测}

与 PhysFlow 侧重于从视觉观测优化材料参数不同，NGFF（2026.02）~\cite{2602_00148}（Neural Gaussian Force Field）从力场预测的角度重新审视 3DGS 物理仿真——不依赖 MPM 等传统数值求解器推断材料属性，而是直接学习一个端到端的神经力场。

\textbf{动机。}MPM 等数值求解器虽然物理可解释，但计算开销受限于粒子数量和网格分辨率，且需要预先指定材料类型；NGFF 用数据驱动的神经算子替代数值仿真，在牺牲部分物理解释性的前提下换取数量级的加速。其力场预测器学习全局变换力场和局部应力场，通过神经算子在关系图上建模物体间物理接触，速度比先前高斯仿真器快两个数量级。这一设计的动机在于：PhysFlow 的可微分 MPM 框架虽然物理可解释性强，但计算开销受限于粒子数量与网格分辨率，且材料类型需要预先指定；NGFF 用数据驱动的神经算子替代数值仿真，在牺牲部分物理解释性的前提下换取了数量级的加速和更强的泛化能力。同时提出 GSCollision 数据集，包含超过 64 万渲染物理视频（约 4TB），涵盖多种材料和复杂交互场景。

\paragraph{Learning from Real-world Videos：真实视频数据驱动的碰撞动力学}

在 PhysFlow 和 NGFF 分别从材料参数优化和力场预测推进 4D 物理仿真的基础上，后续工作进一步拓展了训练数据来源和感知模态。与 PhysFlow 依赖合成参考视频和光流监督不同，Learning from Real-world Videos（2026.05）~\cite{2605_23845}首次实现从真实世界无标注视频通过渲染监督学习多物体碰撞动力学，直接在密集高斯基元（无启发式子采样）上操作，并贡献约 500 视频的真实数据集——这解决了 PhysFlow 框架中视频扩散模型生成的参考运动可能偏离真实物理分布的问题。

\paragraph{PEGS：事件相机增强的高速运动感知}

PEGS（2025.11，Physics-Enhanced Gaussian Splatting）~\cite{2511_17116}则从感知模态角度出发：PhysFlow 的光流监督仅利用标准 RGB 帧，无法捕获高速运动中的瞬间状态；PEGS 融合物理约束、事件相机和 3DGS，实现大时空尺度刚体运动的高质量去模糊重建，贡献首个针对自然快速刚体运动的 RGB-Event 配对数据集。

\textbf{动机。}标准 RGB 相机在高速运动场景中产生严重运动模糊，导致 3DGS 重建质量急剧下降；事件相机以微秒级分辨率捕获亮度变化，为去模糊重建提供了互补的时序信息。

\subsection{声音与多物理场}

物理仿真不仅限于视觉——声音、热量、电磁等多物理场（multi-physics）现象同样可以用高斯基元表示和模拟。本节的方法将 3DGS 从纯视觉表示扩展到多模态物理场。下面先讨论声音场重建与合成，再介绍火焰、降雨、风等其他物理现象的建模方法。

\subsubsection{声音频高斯的显式重建}

本路线介绍将3DGS从纯视觉表示扩展到声学域和多物理场的方法，包括声音场重建、碰撞声合成、火焰/降雨/风等物理现象的建模与仿真。

\paragraph{声音场重建与合成}

AudioGS（2026.04）~\cite{2604_08967}将 3DGS 从视觉域迁移到声学域——用显式"音频高斯"替代隐式神经声场表示，实现不依赖视觉先验的声场重建与双耳合成。在此之前，新视角声学合成（NVAS，Novel View Acoustic Synthesis）的主流做法依赖视觉引导：ViGAS 用跨模态 Transformer 从视频帧提取空间线索，AV-NeRF 用 NeRF 联合渲染 RGB/深度图来指导声学合成，AV-Cloud 则用 SfM 稀疏点云作为声学锚点。然而，这些方法存在两个问题：其一，依赖视觉先验，在纹理缺失、相机标定不准或视觉数据稀疏时性能显著退化。其二，仅重用源视角相位而忽略相位建模，导致耳间时间差（ITD，Interaural Time Difference）线索不准，空间感失真。因此 AudioGS 选择完全抛弃视觉依赖，直接从稀疏声学观测显式重建 3D 声场。

直觉上，AudioGS 将声场编码为一组"音频高斯"（Audio Gaussians），每个音频高斯对应 STFT（Short-Time Fourier Transform，短时傅里叶变换）频谱图上的一个时频 bin（time-frequency bin），配备可学习的 3D 位置、双球谐系数和距离衰减系数。这一设计利用了语音信号在时频域的近似 W-disjoint 正交性——每个时频 bin 的能量通常由单一声源主导——从而可以将每个 bin 独立映射为一个空间基元。双 SH 参数化将方向性能量与空间双耳线索解耦：单耳 SH 系数 $c^{\text{mono}}_{f,t}$ 产生非负的共享能量模式，差分 SH 系数 $c^{\text{diff}}_{f,t}$ 编码左右耳差分线索。渲染时，对目标监听位置 $p$ 和朝向 $R_{\text{cam}}$，AudioGS 依次执行三步操作。（1）方向性调制：通过球谐基函数 $Y(d_{f,t})$ 将 SH 系数投影到从音频高斯指向监听者的方向 $d_{f,t}$ 上，得到单耳掩码 $M(f,t) = 2\,\text{sigmoid}(c^{\text{mono}}_{f,t} \cdot Y(d_{f,t}))$ 和差分掩码 $D(f,t)$。（2）距离衰减：建模频谱幅度随距离的物理衰减。

\begin{equation}
G(f,t) = \left( \frac{\|p_{\text{ref}} - x_{f,t}\|_2 + \epsilon}{\|p - x_{f,t}\|_2 + \epsilon} \right)^{\!\alpha_{f,t}}
\end{equation}

其中 $\alpha_{f,t} > 0$ 是可学习的衰减指数，$\epsilon$ 为数值稳定项。该公式允许模型自适应地从混响场（$\alpha \approx 0$）过渡到直达声衰减（$\alpha \gtrsim 1$）。（3）相位校正：基于刚球模型计算物理 ITD 幅值，并通过有界残差因子进行可学习校正，弥补传统方法直接复用源相位的缺陷。最终将幅度与校正相位组合，经逆 STFT 输出双耳波形。

实验在 Replay-NVAS 数据集上进行，音频重采样至 16\,kHz，STFT 使用 512 点 FFT、400 窗长、160 帧移和 Hamming 窗。训练 60 个 epoch，所有 SH 系数初始化为 0，衰减系数初始化为 1。与视觉引导方法相比，AudioGS 在不使用任何视觉数据的前提下，MAG 误差降低 14\%+，ENV 降低约 6\%，感知质量指标 DPAM 降低约 25\%，耳间能量比误差（LRE）也达到竞争性水平。消融实验证实三个组件缺一不可：移除距离衰减导致空间精度退化最严重，移除 SH 导致重建质量显著下降，移除相位校正则使 LRE 指标明显劣化。MUSHRA 听觉测试（12 名受试者）也表明 AudioGS 的合成音频最接近真实录音。

与 AudioGS 聚焦于声场的空间重建与双耳合成不同，SonicGauss（2025.07）~\cite{2507_19835}关注的是碰撞声合成——不是重建已有声音的空间分布，而是从 3DGS 几何表示主动生成物理碰撞声。AudioGS 的音频高斯编码的是环境中已存在的声场信息，而 SonicGauss 需要从高斯椭球体的几何特征推断材质特性（金属、木质、玻璃等）和空间声学相关性，再通过扩散声音合成模型生成对应的声音。这一差异的设计动机在于：AudioGS 解决的是"听者在不同位置听到什么"的空间声学重建问题，而实际物理仿真场景中更基础的需求是"物体碰撞时发出什么声音"——只有先生成碰撞声，才能进一步讨论其空间传播。SonicGauss 整合 PointTransformer 特征提取器与扩散声音合成模型，支持空间变化的声音响应——同一物体在不同位置被敲击产生不同的回声效果。

声音之外，3DGS 的显式点云表示和可微分渲染管线同样适用于火焰、降雨、风等其他物理现象的建模与仿真。

\paragraph{FlameGS：高斯基元火焰光场重建}

FlameGS（2024.12）~\cite{2412_19841}修改了火焰化学发光断层成像（CTC）的三维表示环节——用 3D 高斯基元替代传统体素网格，实现火焰光场的高效重建。传统方法采用 ART（代数重建技术）：将火焰空间离散为体素网格（如 $120^3 = 1{,}728{,}000$ 个体素），通过 2D 投影方程逐体素迭代更新发光强度。然而体素表示存在两个瓶颈：网格分辨率固定，重建区域受预定义边界严格限制；体素数量庞大导致计算和内存消耗极高。FlameGS 的核心思路是：火焰本身是连续的发光体，可以自然地用一组可微的 3D 高斯函数来近似——每个高斯由位置 $\boldsymbol{\mu}$、各向异性协方差 $\boldsymbol{\Sigma} = \mathbf{R}\mathbf{S}\mathbf{S}^\top\mathbf{R}^\top$（$\mathbf{S}$ 为缩放向量，$\mathbf{R}$ 为旋转四元数）和发光强度 $\sigma$ 三个参数描述，整体仅需约 13{,}000 个可训练参数。直觉上，每个高斯对应火焰中一个"发光团"，其位置、形状和亮度共同刻画火焰的三维光场分布。训练时利用 10 台同步相机（$800 \times 1024$ 分辨率）捕获的 2D 投影图像作为监督，最小化 L1 损失与 D-SSIM 损失的加权和来联合优化高斯参数与相机位姿；初始化阶段则通过光线追踪从多视角反向投影，只在真实火焰区域生成稀疏初始点云，有效缓解稀疏视角下的过拟合。相比传统 ART 算法，FlameGS 计算时间减少约 34 倍，内存减少约 10 倍，同时达到 0.96 的 SSIM 和 39.05 dB 的 PSNR。

\paragraph{FieryGS：MLLM 驱动的在线燃烧模拟}

与 FlameGS 聚焦于\emph{火焰光场的离线重建}（从多视角投影中反推火焰三维光场分布）不同，FieryGS（2026.05）~\cite{2605_00177}从\emph{在线燃烧模拟}角度出发，修改了 3DGS 场景编辑中的物理属性推理与特效生成方式——用 MLLM 驱动的零样本材料推理替代手工标注的可燃性参数，用 Taichi 体积燃烧模拟引擎替代预设视觉特效资产，再通过统一体积渲染器将模拟结果与 3DGS 场景融合。具体而言，框架先将重建的 3DGS 场景分割为语义一致区域，每个区域渲染为 2D 图像后送入 MLLM 推理材料类型、可燃性、热扩散系数和烟雾颜色等燃烧属性；随后 Taichi 引擎基于这些属性执行物理燃烧模拟，处理火焰沿可燃表面传播、烟雾体积扩散和表面碳化等复杂现象；最终统一体积渲染器将燃烧体素与 3DGS 高斯基元在同一管线中渲染，确保火光照明与场景几何的一致性。直觉上，可以理解为三个模块各司其职：MLLM 充当"材料专家"，自动判断场景中每个物体能否燃烧、如何燃烧；物理引擎充当"燃烧模拟器"，负责火焰传播和烟雾扩散的物理计算；渲染器充当"画面合成器"，负责将模拟结果无缝融合到 3DGS 场景中。这是将物理准确燃烧模拟集成到 3DGS 的首个工作，支持火焰强度、气流和点火位置等参数的用户交互控制。

\paragraph{RainyGS：基于物理的实时降雨仿真}

与 FlameGS 和 FieryGS 均围绕\emph{火焰}展开不同，RainyGS（2025.03）~\cite{2503_21442}转向\emph{降雨}这一多相流动场景，修改了传统雨效果合成中的场景交互建模方式——用基于物理的降雨粒子系统与浅水模拟替代简单的 2D 雨滴叠加或视频风格迁移，将其集成到 3DGS 的实时渲染框架中。降雨仿真的物理原理分为两层：空中部分将雨滴建模为球状 3D 高斯基元，沿重力方向排列形成雨丝，通过粒子系统控制雨滴的生成、运动和消亡；地面部分则在二维交错 MAC 网格上执行浅水模拟，处理雨滴落地后的飞溅、积水和涟漪传播，并通过屏幕空间反射（SSR）技术实时计算湿润表面的镜面反射效果。该方法以超过 30fps 的速度生成真实感强、物理准确的动态雨效果，支持从细雨到暴雨的强度灵活控制，在 Mip-NeRF 360 和 Tanks and Temples 等真实场景及大规模驾驶场景上均表现有效。

\paragraph{DiffWind：可微分风-物体交互建模}

与前述方法均针对\emph{单一物理现象}（火焰、降雨）不同，DiffWind（2026.03）~\cite{2603_09668}聚焦于\emph{风-物体交互}这一跨模态耦合问题，修改了动态场景重建中的运动建模方式——用可微分物理模拟驱动的风场-物体联合优化替代纯数据驱动的变形场学习，提出首个物理信息可微分风-物体交互建模框架。其核心创新在于为风和物体分别选择最适合的表示并进行耦合交互：风被建模为规则的网格物理场，每个网格节点存储密度、速度和力等物理量，用 Lattice Boltzmann Method（LBM）模拟其演化以符合流体力学基本定律；物体则继承 3DGS 的粒子表示，每个高斯基元携带外观、材质和运动属性；两者之间的交互通过 Material Point Method（MPM）实现——MPM 将粒子属性映射到背景网格上计算受力，再将更新后的物理量传回粒子，自然桥接了连续场（风）与离散粒子（物体）两种表示。在重建阶段，通过可微分模拟与可微分渲染的端到端反向传播，联合优化时空风场和物体运动参数；在 LBM 物理约束的保障下，优化得到的风场不仅匹配观测运动，还满足流体守恒律。该框架支持前向模拟（在新风条件下预测物体运动）和风力重定向（将重建的风场效果迁移到新场景），并构建了 WD-Objects 风驱动场景数据集用于系统评测。

\paragraph{物理场可视化}

GS-Surrogate（2026.04）~\cite{2604_06358}修改了集合仿真（ensemble simulation）参数空间探索的空间表示与变化建模方式——用\textbf{可变形显式高斯场}替代传统的\textbf{端到端 2D 图像代理}或\textbf{单统一隐式 NeRF 场}，首次实现同时调整仿真参数和可视化设置的实时交互探索。

\textbf{Baseline 做法与问题。}集合仿真是科学计算中的核心需求：研究者对同一物理系统在多种参数设置下运行仿真（如宇宙学中调整物质密度 $\Omega_m$、重子密度 $\Omega_b$ 和哈勃常数 $h$），以理解参数对结果的影响。传统做法有两类：(1) 保存高分辨率仿真输出需要海量存储和 I/O，in-situ 渲染虽减少存储但固定渲染配置，牺牲后分析灵活性；(2) 可视化代理模型（如 InSituNet、ViSNeRF）直接从参数合成 2D 图像或用 NeRF 建模场景，但前者缺乏几何一致性（无法在不同视角保持结构一致），后者因密集射线采样导致渲染速度慢，且将所有参数驱动的变化编码在单一隐式场中——仿真参数变化与可视化设置变化相互纠缠，学习传递函数的变化可能干扰仿真参数的几何差异。

\textbf{核心思路：规范场 + 两阶段变形。}GS-Surrogate 的核心洞察是：不同仿真成员共享大量几何结构，差异可由高斯基元的位置偏移和属性变化捕捉。因此，先用一组共享的\textbf{规范高斯场}（canonical Gaussian field）$\mathcal{G}$ 捕获所有成员的基础几何，再通过条件变形网络将其适配到特定参数组合。关键是将变形分解为两个\textbf{顺序阶段}，将仿真参数驱动的变化与可视化参数驱动的变化显式分离：

\begin{equation}
\mathcal{G}' = F_{\text{sim}}(\mathcal{G}, P_{\text{sim}}), \quad \mathcal{G}'' = F_{\text{vis}}(\mathcal{G}', P_{\text{vis}}),
\end{equation}

其中 $F_{\text{sim}}$ 以每个规范高斯的位置 $\mu_i$ 和仿真参数向量 $P_{\text{sim}}$ 为输入，预测五类属性偏移量：位置偏移 $\Delta\mu$ 捕捉结构位置变化（如风应力增大使洋流结构扩张），缩放偏移 $\Delta s$ 和旋转偏移 $\Delta q$ 捕捉局部几何变化，不透明度偏移 $\Delta o$ 建模拓扑变化（如宇宙学仿真中某些纤维结构在特定参数下消失），颜色偏移 $\Delta c$ 捕捉标量场值变化导致的外观差异。$F_{\text{vis}}$ 在仿真变形结果上进一步调整以适配可视化任务（如传递函数编辑或等值面提取）。这种模块化设计使同一套规范高斯可通过不同的可视化变形复用于多种分析任务，无需重新训练。

\textbf{设计动机。}这一分解源于两个考量：(1) 规范场与可视化参数解耦，使得新增可视化设置时无需从头重训；(2) 显式高斯基元可由可微分光栅化实时渲染，避免了 NeRF 的逐射线查询开销。在 Nyx 宇宙学（暗物质密度场，$512^3$ 体素）、MPAS-Ocean 海洋环流、XCompact3D 湍流和 CloverLeaf3D 冲击波四个数据集上，GS-Surrogate 在体渲染和等值面渲染两种任务中均实现了实时参数空间探索，渲染质量优于 ViSNeRF（PSNR 提升 1--3 dB），推理速度达到交互级别。

\subsection{其他物理方法}

除了上述按材料类型划分的方法，还有一些工作从不同角度推进 3DGS 物理仿真的边界。

\paragraph{哈密顿力学与能量守恒}

\subsubsection{哈密顿力学与能量守恒约束}

本路线介绍利用哈密顿力学、拓扑推断等理论工具为3DGS物理仿真注入结构性约束的方法，以及VR交互与系统集成等应用方向。

NeHaD（2025.12）~\cite{2512_10424}（Neural Hamiltonian Deformation）修改了动态高斯变形场的解码环节——用哈密顿神经网络（Hamiltonian Neural Network, HNN）替代标准 MLP 解码器，使变形场内嵌物理约束。下面先介绍该方法的动机和核心设计。

\textbf{动机。}标准 4DGS 的做法是：hex-plane 编码器提取时空特征后，由纯数据驱动的 MLP 解码器直接预测高斯属性偏移量（位置、缩放、旋转）。这种 MLP 缺乏物理归纳偏置，常产生物理不合理的运动——高斯突变式出现或消失、轨迹不连续、违反能量守恒。NeHaD 的核心思路是：高斯基元天然存在于相空间（位置-动量空间）中，其演化应遵循哈密顿力学，因此用 HNN 替代 MLP 以注入能量守恒的结构性约束。

\textbf{哈密顿变形场。}直觉上，哈密顿力学要求运动的速度场可分解为两种成分：保守场（能量守恒，沿梯度方向推动运动）和螺线场（体积守恒，产生旋转运动）。NeHaD 不直接学习向量场（高维易模式坍缩且难保证能量约束），而是学习两个标量势函数 $F_1$ 和 $F_2$，通过自动微分生成具有物理保证的向量场：

\begin{equation}
\mathbf{v}_c = \nabla_{h_i} F_1(h_i), \quad \mathbf{v}_s = \nabla_{h_i} F_2(h_i) \, M^\top,
\end{equation}

即保守场取标量势的梯度，螺线场取标量势的梯度再乘以辛置换张量 $M$（辛结构是哈密顿力学的核心几何性质，保证相空间体积不变性）。$\nabla \times \mathbf{v}_c = 0$ 保证保守性，$\nabla \cdot \mathbf{v}_s = 0$ 保证螺线（散度自由）性。变形位移由轻量属性特定适配器（线性层）处理：$\Delta\mu_i = A_\mu \mathbf{v}$, $\Delta s_i = A_s \mathbf{v}$, $\Delta r_i = A_r \mathbf{v}$。

\textbf{时间积分。}为保证时间演化的数值稳定性，NeHaD 采用二阶辛积分器（Position Verlet）更新高斯位置：

\begin{equation}
\mu_i = \mu_i + \Delta t \, v_i + \frac{(\Delta t)^2}{2} F_i,
\end{equation}

相比欧拉法，辛积分器在长时间演化中保持能量守恒，避免数值漂移。

\textbf{玻尔兹曼平衡分解。}此外，NeHaD 引入玻尔兹曼平衡分解（Boltzmann Equilibrium Decomposition, BED）自适应分离静态和动态高斯：仅远离平衡态的高斯参与变形，基于 Boltzmann 分布构建软掩码
\begin{equation}
M_{\text{pos}}^{(i)} = (1-\epsilon) \exp(-\beta E_{\text{st}}^{(i)}) + \epsilon,
\end{equation}

即能量偏差 $E_{\text{st}}^{(i)}$ 越大（越偏离平衡态）的高斯越可能被激活，$\epsilon$ 保证最小激活率防止梯度消失，使计算资源集中在真正运动的高斯上。

前述方法聚焦于动力学建模——让高斯"动起来"。另一个方向则关注结构推断——从视觉观测中"看透"物体内部。下面介绍拓扑推断的代表性工作。

\paragraph{拓扑推断与结构发现}

TopoGaussian（2025.03）~\cite{2503_12343}修改了不透明物体内部拓扑推断的表示与仿真管线——用\textbf{全粒子化可微仿真器}替代传统的\textbf{网格依赖管线}，首次从易获取的照片和视频推断不透明物体的内部拓扑结构，无需 CT/MRI 等昂贵设备。

\textbf{Baseline 做法与问题。}推断不透明物体的内部结构（如平衡鹰玩具的内部空腔、古代欹器的非对称内壁）是工程和制造领域的经典难题。传统方法需要 CT 扫描或 MRI 成像，设备昂贵且操作复杂。近年来的 3D 重建方法（NeRF、3DGS）仅恢复外表面几何，不提供内部拓扑。若使用网格方法（如 PGSR、Gaussian Surfels）从点云提取表面网格再进行拓扑优化，则面临网格质量依赖——低质量网格需要繁琐的修补和填充，且网格提取本身就是计算瓶颈。

\textbf{核心思路：全粒子管线。}TopoGaussian 的核心洞察是：利用 3DGS 的点云输出，将建模、仿真和优化三个阶段统一在同一粒子表示上，完全跳过网格提取步骤。具体流程分为四步：(1) \emph{体积形状生成}——用 3DGS 从多视角图像重建外表面点云，再通过体素化填充内部粒子（遵循 PhysGaussian 的内部填充策略），得到密集体积点云 $\{(\mathbf{x}_i, V_i)\}_{i=1}^N$；(2) \emph{拓扑表示}——在每个粒子上附着拓扑参数，提供三种灵活的表示选项：(a) 粒子表示——为每个粒子分配一个决策变量，通过 sigmoid 映射为二元指示函数（$r(\mathbf{x})=1$ 表示实心，$r(\mathbf{x})=0$ 表示空心），最为灵活；(b) 神经隐式表面——用 DeepSDF 网络编码内部边界的符号距离函数 $s(\mathbf{x})$，将 SDF 通过 sigmoid 映射为拓扑指示，能编码几何先验；(c) 二次曲面——用仅 10 个参数的二次函数 $s(\mathbf{x}) = \mathbf{x}^\top A \mathbf{x} + \mathbf{b}^\top \mathbf{x} + c$ 参数化拓扑边界，适用于刚体场景（已证明对刚体拓扑优化充分）；(3) \emph{可微分仿真}——基于粒子系统构建仿真器，同时支持刚体（Newton-Euler 方程，RK2 积分）和软体（Saint Venant--Kirchhoff 超弹性模型，Leapfrog/隐式 Euler 积分），通过核函数估计变形梯度 $F_i$，还集成了气动执行器模型（气压力 $E_a = \sum_i p |\det F_i| V_i$）和罚函数碰撞模型；(4) \emph{梯度优化}——以仿真运动与参考视频之间的 $L_2$ 损失为目标，通过可微分仿真器反传梯度，用 L-BFGS-B 或 Adam 优化拓扑参数。

\textbf{设计动机。}全粒子管线避免了网格提取这一瓶颈步骤。实验表明，相比网格方法 PGSR，TopoGaussian 在体积形状生成阶段平均快 5.26 倍，且重建质量（表面法向量拉普拉斯光滑度）平均优 2.33 倍。在刚体平衡（悬挂鹰、不倒翁）、软体变形（气动软手）和真实物体（平衡鹰、欹器）等任务上验证了有效性，其中一个优化结果已成功 3D 打印并验证物理行为。

DynamicTree（2025.10）~\cite{2510_22213}是首个为 3DGS 重建的真实树木生成交互式 3D 运动动画的框架。其核心表示创新是引入稀疏体素谱（sparse voxel spectrum）来紧凑编码树的运动。

\textbf{动机。}树木的分支结构复杂且数量庞大，逐步物理仿真计算代价极高；稀疏体素谱将树的振动分解为频率模态，动画生成时无需逐步仿真，以前馈方式直接合成各模态叠加结果。直觉上，就像将树的振动分解为频率模态，类似音乐中的泛音：低频模态对应树干的整体摆动，高频模态对应细枝的颤动。通过模态分析（modal analysis），树的复杂空间运动被投影到一组正交频率基上，每个基函数描述一种典型的振动模式。

动画生成时，无需逐步运行物理仿真，而是以前馈方式直接合成各模态的叠加结果。当施加外力（如风力）时，稀疏体素谱还可作为快速模态分析的基底，实现实时交互式响应——树枝随风摆动的方向和幅度由外力频率与各模态固有频率的共振关系决定。该管线先将 3DGS 重建的树转化为网格，利用稀疏体素谱生成网格运动，再将高斯基元绑定到变形网格上完成渲染。为支撑模型训练，DynamicTree 构建了首个大规模 4D 树数据集 4DTree（8786 棵树，每棵 100 帧序列），相比逐步物理仿真加速 100 倍以上。

\paragraph{VR/交互与系统集成}

LIVE-GS（2024.12）~\cite{2412_09176}修改了 VR 物理资产创作中的参数配置环节——用\textbf{LLM 自动化参数预测}替代\textbf{手动逐参数调优}，首次实现从静态 3DGS 资产到可交互 VR（Virtual Reality，虚拟现实）物理场景的全自动转换。

\textbf{Baseline 做法与问题。}VR-GS 等现有 3DGS 物理仿真系统需要用户手动设定每个物体的物理参数（质量、刚度、阻尼等）。作者访谈 7 位物理仿真专家后发现：手动调参既依赖材料科学专业知识又极其耗时，且随场景中交互物体数量增加呈非线性增长——物体间的耦合效应迫使创作者反复修改已完成的参数配置。基于视频扩散模型的自动估计方法（如 PhysDreamer、DreamPhysics）虽可避免手动调参，但计算代价高昂、收敛慢，不适合 VR 资产的即时创作需求。

\textbf{核心思路：LLM 分析 + 实时仿真。}LIVE-GS 的核心洞察是：VR 交互不需要物理精度，只需"看起来对"——而 LLM（GPT-4o）具备足够的视觉-物理常识来推断"看起来对"的参数。系统分两阶段工作。（1）\emph{LLM 驱动的物理分析}：对输入图像，GPT-4o 通过精心设计的提示链（prompt engineering）和函数工具集成（function tool），依次推理物体类别、材料属性和质量，输出结构化 JSON 文件直接对接物理引擎，整个分析过程约 10 秒。（2）\emph{实时物理仿真}：选择 XPBD（Extended Position-Based Dynamics，一种基于位置约束的实时物理求解器，相比 MPM 计算量小两个数量级）作为仿真引擎。对于弹性体，通过 alpha shape 算法从高斯基元重建中间网格表示，在其上生成不可见物理粒子（体素采样密度 $d=15$，内部每个体素中心放一个粒子防塌缩，边界在网格顶点放置粒子保几何）。然后为每个高斯基元找到最近的 $m=4$ 个物理粒子，通过距离反比加权的球面线性插值（SLERP，Spherical Linear Interpolation）将物理粒子的运动传递到高斯：

\begin{equation}
\mathbf{p}^k = \mathbf{p}_0^k + \underbrace{\sum_{i=1}^{m} w_i \Delta \mathbf{p}_i}_{\text{平移}} + \underbrace{\sum_{i=1}^{m} w_i (q_i \, \mathbf{r}_i \, q_i^{-1} - \mathbf{r}_i)}_{\text{旋转引起的平移}}, \quad q^k = \text{SLERP}_{w_i}(q_i) \cdot q_0^k,
\end{equation}

其中 $w_i$ 为初始化时计算的逆距离权重，$\Delta \mathbf{p}_i$ 和 $q_i$ 为第 $i$ 个物理粒子相对初始状态的平移偏差和旋转偏差，$\mathbf{r}_i$ 为高斯基元到物理粒子的初始距离向量。对于刚体，仿真直接在网格上操作，每个高斯基元绑定到最近网格的局部坐标系；当外力超过阈值时通过 Voronoi 分割实现断裂效果。

\textbf{设计动机与验证。}这一设计平衡了交互质量与预测效率：LLM 提供``近似但可信''的参数（在 ABO 数据集 11 个物体上质量预测的 MAPE 为 15.0\%），XPBD 提供实时仿真反馈（渲染帧率 30 fps）。16 人用户研究（VBench 评估）表明，LLM 自动配置的弹性参数在视觉质量、运动平滑度、美学质量和视频-文本一致性四个维度上均达到满意水平，且创作时间从手动调参的数十分钟缩短至约 10 秒。

在综述方面，"Grounding Creativity in Physics"（2025.02）~\cite{2502_07007}系统综述了物理先验在 AI 生成内容中的应用，按表示类型（视觉、NeRF、3DGS）分类方法，揭示了生成模型与物理真实感之间的差距。该综述的核心贡献是首次系统梳理了从纯视觉生成到物理感知生成的技术演进路线，并指出了当前方法在物理精度验证标准上的缺失。

\begin{takeaway}
物理约束与仿真的核心挑战是为没有质量、速度和内部应力的"彩色气泡"赋予物理属性。三条技术路线——粒子绑定（高斯基元绑定到物理粒子）、连续介质力学（MPM 将高斯视为物质点）、以及数据驱动（视频/扩散先验推断物理参数）——分别以不同的精度-效率权衡桥接视觉表示与物理引擎。LIVE-GS 用 LLM 自动预测参数将 VR 物理资产创作从数十分钟缩短至 10 秒。
\end{takeaway}

\subsection{近期方法脉络：场景级物理交互、物理可用重建、物理约束动态高斯等}

这一组近期工作延续 3DGS、Gaussian splatting、4D Gaussian Splatting，主要把改动集中在增量改进、应用、扩展等环节。下面只保留相对既有路线的变化，便于和前文方法直接对照。

\paragraph{Scene-Level Heterogeneous Physics Simulation with 3D Gaussian Splats}
Scene-Level Heterogeneous Physics Simulation with 3D Gaussian Splats（2026.06）~\cite{2606_21753} 延续 3DGS 在“场景级物理交互”上的基本设定，主要改变是：Scene-Level Heterogeneous Physics Simulation with 3D Gaussian Splats：把3DGS资产接入复杂静态碰撞几何和异质物理对象。3DGS 的显式高斯表示和可微渲染能力保持为基础。
其设计动机是：突破以往physics-for-3DGS仅能在理想平面或孤立对象上演示的局限，使3DGS资产可与复杂碰撞几何和多类物理资产交互。
具体而言，方法为3DGS资产建立物理可理解代理或耦合结构，并在生产级物理引擎中进行场景级仿真。
适用边界是：3DGS到物理代理的转换可能损失细节；复杂材料和可变形物体仍难。

\paragraph{Real-Time Physics Simulation with Dynamic Mesh-Gaussian Reconstructions}
Real-Time Physics Simulation with Dynamic Mesh-Gaussian Reconstructions（2026.05）~\cite{2606_00444} 延续 Gaussian splatting 在“物理可用重建”上的基本设定，主要改变是：固定拓扑mesh用于物理碰撞，Gaussian splatting用于外观渲染的双表示框架。3DGS 的显式高斯表示和可微渲染能力保持为基础。
其设计动机是：分析动态Mesh-Gaussian重建中变拓扑几何难以接入物理引擎的问题，并提出双表示方案通过运行时顶点缓冲更新加速仿真集成。
具体而言，系统将重建结果转换为固定拓扑mesh以支持碰撞检测和物理仿真，同时保留高斯表示进行照片级渲染，并比较时序对应跟踪和模板投影等拓扑转换策略。
适用边界是：拓扑转换造成65-80\%几何退化；物理可用性与视觉重建质量之间仍有明显折衷。

\paragraph{PersistGS}
PersistGS（2026.06）~\cite{2606_03479} 延续 4D Gaussian Splatting 在“物理约束动态高斯”上的基本设定，主要改变是：PersistGS：将可微刚体仿真与4D Gaussian Splatting耦合，估计摩擦和速度以维持遮挡期间物体轨迹。3DGS 的显式高斯表示和可微渲染能力保持为基础。
其设计动机是：提出把物体级高斯、碰撞mesh和可微刚体仿真结合起来，在观测消失时用物理轨迹替代纯视觉梯度。
具体而言，方法将场景分解为对象高斯和碰撞网格，从遮挡前轨迹估计速度和摩擦，经可微仿真预测SE(3)轨迹并驱动物体高斯在遮挡期间演化。
适用边界是：假设对象刚体且碰撞模型可得；复杂非刚体和多物体接触会更困难。

\paragraph{Occlusion-Robust Multi-Object Decoupling for Physics-Based Robotic Interaction}
Occlusion-Robust Multi-Object Decoupling for Physics-Based Robotic Interaction（2026.06）~\cite{2606_29303} 延续 3DGS 在“Gaussian reconstruction for physical interaction”上的基本设定，主要改变是：以 3DGS 为基础表示，结合 SAM2 分割场、SDS 和 MPM 物理仿真。3DGS 的显式高斯表示和可微渲染能力保持为基础。
其设计动机是：将多物体解耦表述为稀疏视图重建问题，减少遮挡和有限视角造成的对象耦合。
具体而言，先以 3D Gaussian Splatting 获取基础重建，再用 SAM2 训练的 segmentation field 得到粗实例分区，并通过联合 SDS 重建碎片几何，最后用于 Material Point Method 仿真。
适用边界是：SDS 和分割场可能引入生成偏差；物理仿真可用性取决于几何完整度和材料参数估计。

\paragraph{MeGAS}
MeGAS（2026.06）~\cite{2606_23455} 延续 3DGS 在“热机械与相变高斯”上的基本设定，主要改变是：MeGAS：Thermomechanical Dynamic Gaussian Splatting，将相变、温度和机械动态耦合进高斯表示。3DGS 的显式高斯表示和可微渲染能力保持为基础。
其设计动机是：突破现有physics-for-3DGS只关注机械运动的局限，引入不可见但关键的温度物理因素。
具体而言，方法构建热机械动态Gaussian Splatting表示，用温度场和相变状态驱动高斯形变、材质和运动变化。
适用边界是：热物理参数估计和相变模型复杂；真实材料泛化需验证。

\section{近期研究（2026 年 8--11 月）}
\paragraph{从视觉和少量录音估计物体的模态声场}
2026.08~\cite{2608_05145}。这项工作用多视图图像提供物体空间形态的线索，用少量已知接触产生的录音估计模态声学参数。视觉高斯是约束声音合成的空间表示，预测的声音响应还依赖物体的模态参数与敲击位置。

论文评价未见接触位置的撞击声，并展示接触定位等下游任务。少样本指已记录的撞击声较少，方法仍需要对同一物体的视觉观察与物理响应，证据主要对应所测撞击设置。

\paragraph{WildFireGS：在带植被语义的森林高斯上模拟火灾}
2026.08~\cite{2608_11100}。WildFireGS 将航空影像重建的高斯扩展成带植被类别和燃料属性的场景，经过尺度、重力和场景结构预处理后，用粒子模型在高斯表示上模拟热传递与燃料消耗。高斯同时承载可视外观和参与仿真的空间属性。

论文在合成林地控制变量并在真实航空重建场景演示，评价传播趋势与火隔离带和生物量变化。由航拍不完整观测推断的植被结构和燃料参数会影响结果，合成控制实验不等于现场火势预测精度。

\paragraph{LaGSplat：用拉格朗日潜状态驱动动态高斯}
2026.08~\cite{2608_16324}。常规动态高斯以时间索引场景，难以直接响应新激励。LaGSplat 用同一个低维潜向量充当耗散拉格朗日系统的广义坐标和高斯解码器的控制变量。模型从单目视频学习潜状态及其动力学，使原语随物理状态变化而运动。

论文在多个真实系统评价记录运动和片段外外推，并比较实物软体机器人的潜在动力学。视频并不直接监督潜变量，实验展示所测系统的运动与响应，不代表在任意系统中已从单目视频辨识唯一物理参数。

\paragraph{从单目动态视频对齐形变外观与隐式本构模型}
2026.08~\cite{2608_22102}。这项工作结合低秩调整、秩深度几何锚点与本构先验正则，使模拟材料的形变和单目渲染互相约束。前者稳定视觉重建，后者利用现有本构先验限制模拟行为。高斯是可微视觉表示，材料规律由隐式模型学习。

在包含六类材料的合成实验中，作者相对 NeuMA 报告粒子位置 Chamfer 距离平均降低 48\%。在 SpringGaus 的真实数据上另用单目监督验证，采用的评价与具有粒子真值的合成实验不同，不能从合成误差直接推断真实材料属性恢复的精度。

\paragraph{LagrangeGS：从预训练高斯轨迹蒸馏非保守动力学}
2026.08~\cite{2608_22773}。LagrangeGS 先从预训练动态高斯提取每粒子的轨迹，再按对象蒸馏势能、初始速度和非保守力，并用局部刚性对齐避免独立粒子解体。推断可改变学习到的势场以观察反事实运动，或通过负时间步向后积分。

论文在 Dynamic Object 和 Dynamic Indoor Scene 的保留视图与时间窗口评价。前向再反向的平均位置循环误差分别为 1.3\% 和 0.19\%，是数值一致性测量。方法依赖预训练轨迹，解耦粒子假设也限制了碰撞与接触交互。

\paragraph{KnockGS：用一次有控制的响应校准物理高斯}
2026.08~\cite{2608_27365}。KnockGS 固定资产、粒子与 MPM 求解器，只校准材料参数的两个尺度。方法从物体专用响应库选取接近已观测激励的局部样本并进行岭回归，得到的参数冻结写回同一仿真器，再预测未参与拟合的另一组激励。

实验比较最近邻检索、全局回归与固定默认参数，并测试保留的激励方向和幅度。识别建立在已知激励、同一资产和固定仿真合同上，不能直接泛化为任意物体和无控制观测下的材料识别。

\paragraph{PhysMAS：按部件组织材料与运动再做共享域仿真}
2026.09~\cite{2609_07174}。PhysMAS 接收运动描述、四张已标定图像和已对齐的高斯场景。场景模块保持物体及部件身份，材料模块为各部件给出物理配置，再用共享域 MPM 仿真生成动态高斯，避免为每个场景反复做生成梯度优化。

论文比较六类材料场景的语义对齐、感知的物理合理性及参考视频运动一致性。可选本构配置和空间分辨率限制了部件行为与长时互动的覆盖，遮挡下的部件关联也仍可能不稳定。

\paragraph{从一次扫描合成视觉运动策略训练示范}
2026.09~\cite{2609_21112}。该工作从一次 RGB 视频扫描还原背景、各物体的高斯资产，再与按关节配置渲染的机器人模型共同组成场景。物体转移由刚体运动学控制，抓取闭合等接触过渡由学习模型估计，以渲染观测与机器人动作配对形成模仿学习示范。

实验以任务、对象和扫描的组合作为评价单元，每组仿真试验 100 次、实机 10 次。ManiSkill3 的受控对比与真实工作台实验衡量不同条件；平面布局、刚性物体运动和接触模型限制了可覆盖的操作。

\paragraph{GaussPDE：让标量场在静态高斯图上扩散并驱动显示}
2026.09~\cite{2609_27264}。GaussPDE 从静态高斯选择活跃原语，构建稀疏传播图并建立场景相关的导通关系。标量场在图上演化并调节高斯外观，渲染期间高斯几何和图连通性保持固定。

论文用独立合成网格求解验证扩散及跨表面泄漏，并展示三个场景数据集上的完整流程。传播仍依赖重建几何和提升后的视觉线索，结果不能直接推广到任意运动物体或任意物理方程。

\paragraph{PneuTac：用 MPM 和高斯共同模拟软指与触觉图像}
2026.09~\cite{2609_38418}。PneuTac 结合材料点法与高斯渲染模拟气动软指及视觉触觉传感器。使用视觉数据校准后，动作和感知网络可近似部分仿真过程，而经校准的模拟器还用于扩展少量真实示范。

论文以安装 Digit 触觉传感器的气动软指为主要操作实例，并对其他器件做建模验证。实时替代模型与完整 MPM 的适用条件不同，新装置的材料、气动和传感条件仍需重新校准。

\paragraph{用可微的拉格朗日传递连接高斯与 SPH 变形仿真}
2026.10~\cite{2610_04336}。高斯负责从影像重建固体外观，SPH 则用独立离散的材料点满足力学建模需要。该工作在参考构型中建立可微的拉格朗日传递，将 SPH 粒子的位移映射到视觉高斯，并允许影像误差沿该映射影响逆问题求解。

论文分别评价正向运动仿真与从影像反演固体参数。可微连接使参数可以通过影像监督估计，边界条件、各自离散的合理性和高斯包络质量仍是解释误差与参数可辨识性的必要条件。

\paragraph{以接触、运动和应力场连接视频与材料参数推断}
2026.10~\cite{2610_17591}。这项工作从多视图形变视频先推断接触或力、运动与应力等中间场，再预测本构类别和连续材料参数。这把最终参数估计拆解为可检查的物理证据，便于区分追踪准确与材料参数准确。

评价包括参数、稠密场、同模拟器回放与跨协议迁移。默认设置采用三视角 RGB 与物体级划分，屈服应力仅在发生屈服的样本上评价。DeformFieldBench 以合成为主，真实回放缺少已标定力测量和材料真值，不能作为绝对参数精度的证据。

\paragraph{ReStep：以预测终态约束长程粒子轨迹}
2026.10~\cite{2610_18168}。ReStep 先用图式变分模型从初态与物理参数预测合理终态，再用受这两个端点约束的扩散模型补全中间粒子轨迹。与逐步自回归相比，整个预测窗口共享终态约束，高斯渲染用于视觉条件下的物理应用。

实验在 FlexiFall 的落体、倾倒和垂挂场景评价长程误差，并在 Topple 测试视觉物理参数估计。若预测终态错误，它也会作为强条件传到中间轨迹，产生不可信的运动。

\paragraph{用引力约束的动态高斯重建模拟星系}
2026.11~\cite{2611_20089}。这项工作以自发光高斯表示恒星，让亮星沿独立引力轨道运动，并以共享流近似较暗恒星的运动。从多个时刻的带标定视图优化外观与运动时，引力先验约束高斯长时间演化。

实验覆盖八个模拟星系，讨论旋转、角动量、尺寸与渲染质量。实际望远镜对单个星系通常只有单方向和单历元观测，不能提供评估所需的多相机、多时刻数据，因此这组实验说明的是仿真条件下的可行性。

\paragraph{CiGS：从高斯重建表面求解固体内部温度场}
2026.11~\cite{2611_20283}。CiGS 用已配准的可见光和热图重建高斯表面，再提取带热边界观测的点云，直接求解边界积分方程以推断内部温度。深度斜率过滤剔除跨深度断裂的采样点，内部积分对密度常数项单独计算，整个过程不需要表面或体积网格。

作者比较完整覆盖和单视图采样下的精度及耗时，并使用无矩阵的可微求解。由边界观测推断内部温度依赖重建表面和正确的物理边界设定，单视图结果尤其受表面覆盖限制。


\section*{本章小结}
\addcontentsline{toc}{section}{本章小结}

\begin{futurebox}{本章总结与展望}{ch07_summary}
本章综述了 3DGS 物理仿真与动力学领域的 49 篇论文，涵盖以下核心方向：
\begin{itemize}
\item \textbf{物理仿真集成}：从 MPM 紧耦合（PhysGaussian、i-PhysGaussian、GaussianFluent）到松耦合代理结构（GASP、MaGS、PhysTalk），再到物理属性感知（Sim Anything、PIXIE、PhysGS）
\item \textbf{流体动力学}：Gaussian Fluids 的无网格流体求解器、LagrangianSplats 的无散度核约束、FluidGaussian 的仿真不确定性引导重建、以及 SPH-GS 的光滑粒子流体模拟等方向
\item \textbf{可变形物体与布料仿真}：PhysTwin 的弹簧-质点数字孪生、NeuMA 的物理先验+学习修正、3DGSim 的端到端力场动力学、MonoPhysics 的单目逆物理估计
\item \textbf{声音与多物理场}：AudioGS 的显式声场重建、SonicGauss 的碰撞声合成、FieryGS 的燃烧模拟、RainyGS 的降雨仿真、DiffWind 的风-物体交互建模
\item \textbf{理论创新}：NeHaD 的哈密顿力学变形场、TopoGaussian 的内部拓扑推断、PIDG 的 Cauchy 动量约束
\end{itemize}

当前该领域最突出的开放问题包括：
\begin{itemize}
\item \textbf{物理精度的验证标准}：大多数方法通过视觉质量（PSNR/SSIM）评估物理仿真结果，但视觉真实与物理真实之间可能存在系统性偏差——需要建立基于物理度量的评估基准
\item \textbf{多材料耦合的统一框架}：现有方法通常针对单一材料类型（弹性/流体/断裂），同时处理刚体、弹性体、流体和断裂的统一仿真框架仍需探索
\item \textbf{实时交互的物理仿真}：MPM 等离线仿真方法计算代价高昂，如何在保持物理精度的同时实现实时交互是工程落地的关键瓶颈
\item \textbf{从观察到仿真的全自动管线}：从单目/稀疏视角视频自动推断完整物理参数、内部拓扑和材料属性并生成仿真就绪资产，仍是长期目标
\end{itemize}
\end{futurebox}
